\pdfoutput=1

\documentclass[11pt]{article}

\usepackage[final]{acl}

\usepackage{times}
\usepackage{latexsym}

\usepackage[T1]{fontenc}

\usepackage[utf8]{inputenc}

\usepackage{microtype}

\usepackage{inconsolata}

\usepackage{graphicx}

\usepackage{amsmath}
\usepackage{amsfonts}
\usepackage{amsthm}
\usepackage{algorithmic}
\usepackage{algorithm}
\usepackage{colortbl} %
\usepackage{tikz}
\usepackage{tikz-3dplot}
\usepackage{xcolor}
\usepackage{booktabs}
\usepackage{graphicx}
\usepackage{etoolbox} %
\usepackage{comment}
 \usepackage{multirow} 
\usepackage{pifont}
\usepackage{lipsum}
\usepackage{graphics}
\usepackage{tcolorbox}
\usepackage{enumitem}



\title{Improving Large Language Model Safety with 
\\
Contrastive Representation Learning}


\author{Samuel Simko%
\\
  ETH Zurich \\
  \texttt{ssimko@student.ethz.ch} \\\And
  Mrinmaya Sachan%
  \\
  ETH Zurich \\
  \texttt{msachan@ethz.ch} \\\AND
  Bernhard Sch\"olkopf \\
  MPI for Intelligent Systems \\
  \texttt{bs@tue.mpg.de} \\\And
  Zhijing Jin \\
  MPI \& University of Toronto \\
  \texttt{zjin@cs.toronto.edu} \\
}

\begin{document}
\maketitle
\begin{abstract}
Large Language Models (LLMs) are powerful tools with profound societal impacts,
yet their ability to generate responses to diverse and uncontrolled inputs
leaves them vulnerable to adversarial attacks. While existing defenses often
struggle to generalize across varying attack types, recent advancements in
representation engineering offer promising alternatives. In this work, we
propose a defense framework that formulates model defense as a contrastive
representation learning (CRL) problem. Our method finetunes a model using
a triplet-based loss combined with adversarial hard negative mining to encourage
separation between benign and harmful representations. Our experimental results
across multiple models demonstrate that our approach outperforms prior
representation engineering-based defenses, improving robustness against both
input-space and embedding-space attacks without compromising standard
performance. \footnote{Our code is available at \url{https://github.com/samuelsimko/crl-llm-defense}}
\end{abstract}

\section{Introduction}

In recent years, Large Language Models (LLMs) have proven to be powerful tools for general-purpose language understanding and generation \citep{minaee2024largelanguagemodelssurvey}.
They have had significant impact on software engineering \citep{hou2024largelanguagemodelssoftware}, medicine \citep{zhou2024surveylargelanguagemodels} and natural
sciences \citep{nejjar2024llmsscienceusagecode}.

However, their ability to respond to uncontrolled inputs comes with serious security risks \citep{geiping2024coercingllmsrevealalmost}, as they can generate inappropriate, toxic, or harmful text \citep{deshpande2023toxicitychatgptanalyzingpersonaassigned}. To mitigate this risk, various mechanisms have been developed to keep LLMs aligned with human values \citep{stiennon2022learningsummarizehumanfeedback, bai2022traininghelpfulharmlessassistant, rafailov2024directpreferenceoptimizationlanguage}.
Nonetheless, the most commonly used LLM systems often fail to protect against such behaviors \citep{zou2023universaltransferableadversarialattacks, chao2023jailbreaking}.
Developing safe and reliable defenses is therefore essential for minimizing societal risks associated with broad AI deployment.

The current state-of-the-art defenses against LLM jailbreaks focus either on creating effective pipelines around the model \citep{zeng2024autodefensemultiagentllmdefense, yuan2024rigorllmresilientguardrailslarge}, or on finetuning the model weights \citep{yousefpour2025representationbendinglargelanguage, zhang2024negativepreferenceoptimizationcatastrophic}. Among the latter, circuit breakers \citep{zou2024improvingalignmentrobustnesscircuit} are a promising approach, designed to disrupt the model's harmful inner representations, causing it to generate incoherent or nonsensical outputs rather than harmful content even under attacks.

In this work, motivated by the success of circuit breaking, we introduce a new approach to representation space safety engineering by building on concepts from contrastive learning \citep{khosla2021supervisedcontrastivelearning}. Our triplet loss formulation represents a natural extension of circuit breaking, and allows for clearer separation of harmful and benign representations. Additionally, we adopt an improved training strategy based on adversarial hard negative mining \citep{casper2024defendingunforeseenfailuremodes}. Overall, we improve the robustness of representation engineering-based defensive methods by reducing the attack success rate (ASR) of the Llama 3 8B model from 29\% to 5\% against embedding attacks across diverse configurations, and from 14\% to 0\% against the REINFORCE-GCG \citep{geisler2025reinforceadversarialattackslarge} input space attack.

Our main contributions are as follows:

\begin{enumerate}
  \item We propose a novel method for LLM safety based on contrastive representation learning, generalizing and improving upon existing methods
    such as circuit breakers \citep{zou2024improvingalignmentrobustnesscircuit} and RepBend \citep{yousefpour2025representationbendinglargelanguage}.
  \item We introduce a new training strategy based on representation adversarial training \citep{casper2024defendingunforeseenfailuremodes} that improves the sampling procedure of harmful representations, further increasing the robustness of our proposed method against embedding attacks.
  \item We provide experimental validation of our method, demonstrating gains in robustness against both input space and embedding space attacks without sacrificing the model's general capabilities.
\end{enumerate}

\begin{figure*}[ht]
    \centering
    \includegraphics[width=\textwidth]{latex/images/main.pdf} %
    \caption{Comparison of the Triplet defense with the Circuit Breaking defense. Contrary to other adversarial defense methods, circuit breaking aims to break generation at harmful content instead of refusing to answer harmful tasks. It fine-tunes models to keep learned harmless states (or representations) close together while separating newly learned harmful states from their original counterparts, without additional constraints. In contrast, the Triplet defense additionally pulls learned harmful states together and pushes them away from learned harmless states, which increases contrast and robustness to embedding-space attacks.}
    \label{fig:main}
\end{figure*}

\section{Problem formulation}

Large Language Models (LLMs) are deep neural networks, typically based on the Transformer architecture \citep{vaswani2023attention}, that have been trained at a massive scale on human textual data.
The computation of a token $y_t$ given previous tokens $y_{<t}$ can be described as follows:
\newcommand{\emb}{\text{Emb}}
\newcommand{\h}{\textbf{h}}
\newcommand{\hp}{\textbf{h}^\prime}
\newcommand{\out}{\textbf{o}}
\definecolor{darkgreen}{rgb}{0.0, 0.5, 0.0}
\begin{align}
	\h_0 &= \mathrm{Emb}(y_{<t}) \\
	\h_{l} &= T_{l}(\h_{l-1}) + \h_{l-1} \\
	\out &= \mathrm{FFN}(\h_L)
\end{align}
where $\h_0 \in \mathbb{R}^{t-1 \times d}$ is the sequence of input embeddings for the previous tokens, $T_l$ is the transformer block at layer $l$,
$\h_l \in \mathbb{R}^{t - 1 \times d}$ is the output of the $l$-th layer, $\mathrm{FFN}$ is a final feed-forward layer,
and $\out \in \mathbb{R}^{t - 1 \times |\Sigma|}$ represent the output log probabilities over the vocabulary $\Sigma$ for each position. Finally, $y_t$ is sampled from the output distribution $\out$.

For a prompt $x$ and a response $y$, we write $\h_l(x, y)$ as the inner representation of the model at layer $l$.
We consider a prompt and response pair $(x, y)_h$ to be harmful if its content violates the developer rules of the model.
Conversely, a pair $(x, y)_b$ benign does not violate developer rules.

\paragraph{Objective}
Our study specifically targets prompts that encourage illegal, immoral, unethical, or dangerous conduct and actions.
The objective is to minimize the probability that a model outputs a harmful reply $y$ under any token or embedding input $x$.

We use the same proxy objective as the circuit-breaking method of working in the representation level, as while input and output formats can change between different generations of the harmful behavior, the representation space symbolizes the same concept under different formats. We aim to create a new safer model that ``breaks'' when given a harmful prompt, and works as expected when given a benign prompt.

\section{Related work}
\subsection{AI Safety}
\paragraph{Input-Space Jailbreaking Attacks}
A jailbreak is a prompt specifically designed to bypass the model's safety mechanisms and elicit a harmful response,
and can be broadly categorized into token-level jailbreaks which optimize a harmful 
adversarial sequence of tokens appended to the prompt (e.g. Greedy Coordinate Gradient \citep{zou2023universaltransferableadversarialattacks}), or prompt-level jailbreak which
optimize the entire prompt into human-readable jailbreak prompts (e.g Prompt Automatic Iterative Refinement (PAIR) \citep{chao2024jailbreakingblackboxlarge}).

\paragraph{Embedding-Space Jailbreaking Attacks}
Embedding attacks directly manipulate the output of the model's embedding layer to produce a harmful response \citep{schwinn2025softpromptthreatsattacking}. For a prompt $x$, instead of optimizing an adversarial suffix $y \in \Sigma^t$, the attacker can optimize an embedding $e \in \mathbb{R}^{t \times d}$
that is appended to the prompt embedding $\mathrm{Emb}(x)$, to force the model to start with a positive reply. Mathematically, we can write $\h_0 = \mathrm{Emb}(x) \oplus e$ where $\oplus$ is the concatenation operator.
Embedding attacks are more powerful than input space attacks as they operate on a lower level and have access to the entire embedding space of the model.


\paragraph{Jailbreak Defenses}

Jailbreak defenses aim to prevent or mitigate the effects of jailbreak attacks on LLMs.
Current defenses fall into two main categories : Pipeline defenses external to the model, which do not modify the model weights but rather add components such as filters \citep{jain2023baselinedefensesadversarialattacks}, textual transformations \citep{robey2024smoothllmdefendinglargelanguage, yuan2024rigorllmresilientguardrailslarge} or guardrail models \citep{zeng2024autodefensemultiagentllmdefense}, and internal defenses which alter the model itself, by fine-tuning on preference data \citep{kaufmann2024surveyreinforcementlearninghuman} or editing problematic layers \citep{zhao2024defendinglargelanguagemodels}.

More details on jailbreaking attacks and defenses are found in Appendix~\ref{app:details_related}.
In this work, we focus on a subset of internal defenses that manipulate the model's internal representations to increase robustness against jailbreak attacks.

\paragraph{Internal Defenses based on Representation Engineering}

Representation Engineering \citep{zou2025representationengineeringtopdownapproach} focuses on internal representational spaces to understand and improve the behavior of LLM,
and is an alternative to mechanistic interpretability. Circuit breakers \citep{zou2024improvingalignmentrobustnesscircuit} and RepBend \citep{yousefpour2025representationbendinglargelanguage} are the two main
representation engineering-based defenses.
These methods share common principles:

\begin{enumerate}
	\item They manipulate the inner representations $\h_l$ of the model, rather than only the inputs and outputs.
    \item They define training loss functions over these inner representations to guide and optimize model behavior.
    \item They use datasets of prompts and responses labeled as ``benign'' or ``harmful'' to establish ``good'' and ``bad'' inner representations.
\end{enumerate}

\newcommand{\norm}[1]{\left\lVert #1 \right\rVert}
\newcommand{\cossim}{\text{cos\_sim}}
\newcommand{\ReLU}{\text{ReLU}}

The training loss of circuit breaking focuses on preserving benign representations and pushing the new harmful representations far from previous harmful representations using a cosine similarity loss, and is defined as
\begin{align}
	\mathcal{L}_{CB} &= \alpha \cdot \norm{\h_{b} - \h^\prime_b}_2^2 \notag\\
    &+ \beta \cdot \ReLU(\cossim(\h_h, \h^\prime_h)) \label{eq:cb_loss}
\end{align}

where $\h$ is the inner representation of the initial fixed model at layer $l$, $\h^\prime$ is the new inner representation of the model with circuit breaking.
$\alpha$ and $\beta$ are scheduling hyperparameters, $\h_b$ and $\h_h$ represent benign or harmful representations respectively,
and $\cossim(a, b) = \frac{a \cdot b}{\norm{a} \cdot \norm{b}}$ is the cosine similarity between two representations $a$ and $b$.

Circuit breaking is an effective defense against a wide range of input-space attacks,
and decreases the quality of successful attacks.
For instance, while the REINFORCE adversarial attack \citep{geisler2025reinforceadversarialattackslarge} achieves a high attack success rate (ASR) with the HarmBench judge \citep{mazeika2024harmbenchstandardizedevaluationframework}, the resulting responses often display stutter-like behavior in later stages of generation, rendering them mostly unusable.
However, circuit breaking is less effective at defending against embedding-level jailbreaking attacks \citep{schwinn2024revisitingrobustalignmentcircuit}.

Representation Bending (RepBend) \citep{yousefpour2025representationbendinglargelanguage} extends circuit breaking by replacing the cosine similarity-based distance with an L2 distance, and introduces an additional regularization term to enforce similarity among harmful representations. However, training this method is challenging,
as the distance terms can diverge to negative infinity, requiring careful stopping criteria and making it difficult
to maintain a balance of hyperparameters.

\subsection{Contrastive Representation Learning}
The objective of contrastive representation learning (CRL) is to train models to produce a representational space in which
similar (positive) inputs are mapped close to each other, while dissimilar (negative) inputs are mapped far apart.
Rather than solely relying on labeled data, contrastive representation can learn meaningful representations
by leveraging the inherent structure of the data itself. It has achieved notable success in a variety of fields,
such as computer vision \citep{Le_Khac_2020, Schroff_2015}, natural language processing \citep{mikolov2013efficientestimationwordrepresentations, rim2021adversarialtrainingcontrastivelearning}, and multi-modal learning \citep{radford2021learningtransferablevisualmodels}.

The triplet loss \citep{Schroff_2015} is a popular loss function used in contrastive learning, originally developed for image embeddings and face recognition and later adopted for text embeddings \citep{reimers2019sentencebertsentenceembeddingsusing}.

It encourages an anchor data point $a$ to be closer to a positive sample $p$ (similar to the anchor) than to a negative sample $n$  (dissimilar to the anchor) by at least a margin $m$:
\begin{align}
	\mathcal{L}_{T} = \ReLU(d(a, p) - d(a, n) + m)
\end{align}
where $d(., .)$ is a distance metric, typically the Euclidean distance or cosine distance, and $\ReLU(x) = \max(0, x)$ is the rectified linear unit function.
A visualization of the triplet loss objective is shown in Appendix~\ref{app:details_related}.

\section{Method}



\newcommand{\dist}{\text{d}}

\definecolor{darkgreen}{rgb}{0.0, 0.5, 0.0}
\definecolor{darkred}{rgb}{0.6, 0.0, 0.0}

\newcommand{\gcheck}{\textcolor{darkgreen}{\checkmark}}
\newcommand{\rcross}{\textcolor{darkred}{\ding{55}}}

We first describe the desirable properties of our new, more robust representation space.
Our proxy objective is to construct a new representation space $\hp$ that induces the following properties, for an index $i$ and benign and harmful representations $\hp_{b,i}$ and $\hp_{h,i}$:

\begin{enumerate}
    \item $\hp_{b,i}$ and $\h_{b,i}$ should be similar to each other, as the new model is expected
    to maintain similar behavior to the original model for benign use cases. Exact matching is not required for all representations, but top logits of benign behaviors should closely match.
  \item $\hp_{h,i}$ and $\h_{h,i}$ should be dissimilar to each other, because if the general structure of the representation space is preserved, the old harmful representations maintain their harmful nature in the new representation space.
  \item $\hp_{b,i}$ and $\hp_{h,i}$ should be dissimilar, allowing the model to distinguish between benign and harmful representations.
  \item $\hp_{h,i}$ and $\hp_{h,j}$ should be similar to each other, preventing the model from generating fine-grained responses to harmful queries and promoting the generation of uniform replies such as refusals or warnings.
\end{enumerate}

\paragraph{Interpreting the circuit breaking loss}
\label{sec:interpret_cb}
The circuit breaking loss described in Equation~\ref{eq:cb_loss}
can be interpreted as a contrastive loss, similar in spirit to the DrLIM loss \citep{hadsell2006dimensionality},
which is one of the first contrastive losses.
\newcommand{\bX}{\textbf{X}}

For input vectors $\bX_1$ and $\bX_2$ belonging to a class $Y \in \{0, 1\}$ the DrLIM loss is defined as
\begin{align}
\mathcal{L}_{DrLIM} &= (1 - Y) \frac{1}{2} \norm{\bX_1 - \bX_2}_2^2 \notag \\
&\quad+ (Y) \frac{1}{2} \max(0, m - d(\bX_1, \bX_2))
\end{align}
where $d(\bX_1, \bX_2)$ is a distance and $m$ is a margin hyperparameter.

This objective reduces to the circuit breaking objective when using the distance $d(\bX_1, \bX_2) = 1 - \cossim(\bX_1, \bX_2)$, a hard margin $m = 1$, harmfulness labels $Y$, and sampling $\bX_1, \bX_2$ from the original and fine-tuned models.
 
In CRL tasks, the DrLIM loss has been largely supplanted by more effective objectives, notably the Triplet loss and the InfoNCE loss \citep{oord2019representationlearningcontrastivepredictive}, as they are more flexible and induce greater contrasts
between the representations. Motivated by these advances, we use a triplet loss to learn a robust representation space for LLM defense.

\subsection{Our Triplet-Based Loss}
\newcommand{\bp}{\mathbf{p}}
\newcommand{\bn}{\mathbf{n}}

Taking inspiration from the circuit breaking loss function, we propose a general alternative loss function that fits all wanted properties.

Let $\dist_{h,p}$, $\dist_{h,n}$, $\dist_{b,p}$ and $\dist_{b,n}$ be distance functions on representations, and $i$ an index.
We define a harmful triplet loss as:
\begin{align}
\mathcal{L}_{triplet}(h_i) &= \ReLU (\text{d}_{hp}(\hp_{h,i}, \bp_{h,i}) \notag\\
  &\quad- \text{d}_{hn}(\hp_{h,i}, \h_{h,i}) + m_h)
\end{align}


This loss encourages new harmful representations $\hp_{h,i}$ to be distant from the old harmful representations $\h_{h,i}$, and close to some positive $\bp_{h,i}$.
In contrast to circuit breaking and RepBend, our approach focuses on relative rather than absolute distances between representations, as relative distances are more meaningful in embedding spaces.
Conversely, we define a benign triplet loss as:
\begin{align}
  \mathcal{L}_{triplet}(b_i) 
  &= \ReLU(\text{d}_{bp}(\h_{b,i}, \hp_{b,i}) \notag\\
  &\quad- \text{d}_{bn}(\hp_{b,i}, \bn_{b,i}) + m_b)
\end{align}

This loss encourages new benign representations $\hp_{b,i}$ to be close to the old benign representations $\h_{b,i}$, and far from some negative $\bn_{b,i}$.
We write the final, unified triplet loss as a weighted sum of the two triplet losses:
\begin{align}
	\mathcal{L}_{triplet} := \alpha \mathcal{L}_{triplet}(b_i) + \beta \mathcal{L}_{triplet}(h_i)
\end{align}
with hyperparameters $\alpha$ and $\beta$ controlling the importance of the losses.

\paragraph{Relation to Circuit Breakers and RepBend}

\begin{table}[htbp]
\setlength{\tabcolsep}{4pt}
\centering
\begin{tabular}{lccc}
\hline
\textbf{Property} & \textbf{RepBend} & \textbf{CB} & \textbf{Triplet} \\
\hline
$\hp_{b,i} \approx \h_{b,i}$ & \gcheck & \gcheck & \gcheck \\
$\hp_{h,i} \not\approx \h_{h,i}$ & \gcheck & \gcheck & \gcheck \\
$\hp_{b,i} \not\approx \hp_{h,i}$ & \rcross & \rcross & \gcheck \\
$\hp_{h,i} \approx \hp_{h,j}$ & \gcheck & \rcross & \gcheck \\
\hline
\end{tabular}
    \caption{Safety representation engineering methods and their properties. 
    Properties are expressed in terms of similarity ($\approx$) or dissimilarity ($\not\approx$)}
    \label{tab:loss_comparison}
\end{table}

We demonstrate that both the circuit breaking loss and the RepBend loss are simplified special cases of our triplet loss.
The full derivation and proofs are in Appendix~\ref{sec:appendix_proofs}.
Table~\ref{tab:loss_comparison} outlines the key differences between the three losses.
Specifically, the circuit breaking loss lacks mechanisms for separating benign representations from harmful ones, and for clustering the harmful representations. The RepBend loss focuses on clustering harmful representations, but does not explicitly separate benign representations from harmful ones. Our triplet loss formulation generalizes both methods by incorporating these properties and allowing any valid pseudodistances.

\newcommand{\KL}{\mathrm{D}_{KL}}
\newcommand{\dd}{\text{d}}
\newcommand{\reddd}{\textcolor{red}{$\dd_0$}}

\newcommand{\cmark}{\textcolor{green!50!black}{\ding{51}}}
\newcommand{\xmark}{\textcolor{red!50!black}{\ding{55}}}

\paragraph{Choice of the Positive and Negative Samples}
The choice of $\bp_{h,i}$ and $\bn_{b,i}$ is important, as these samples will guide the new representations to new, better directions.
In this work, we use the mean of the new harmful representations as a positive sample for the harmful triplet loss, and as a negative sample for the benign triplet loss,
thereby strengthening the separation between the two classes of representations.
A visualization of our method compared to circuit breakers is shown in Figure~\ref{fig:main}. 
Future work should explore the use of other choices of $\bp_{h,i}$ and $\bn_{b,i}$.

\paragraph{Final Triplet Loss} Let $\alpha$, $\beta$ and $\gamma$ be hyperparameters controlling the importance of the losses.
Let $\dd_{bp}, \dd_{bn}, \dd_{hp}$ and $\dd_{hn}$ be distances chosen by the user.
Let $\h_{b,i}$ and $\h_{h,i}$ be the benign and harmful representations for a batch with $N$ different benign and harmful prompts.
Let $\hat{\hp} = \frac{1}{N} \sum_{h=1}^{N} \hp_{h}$ be the mean of the harmful representations for a batch. Let $\KL$ the Kullback-Leibler divergence on benign model logits between the new and the original model.
Our final triplet loss is defined as:
\begin{align}
  \mathcal{L}_{Triplet} &= \alpha \cdot \frac{1}{N} \sum_{i=1}^{N} \mathcal{L}_{triplet}(b_i) \\
                        &+ \beta \cdot \frac{1}{N} \sum_{i=1}^{N} \mathcal{L}_{triplet}(h_i) \\
    &+ \gamma \cdot \KL(\mathcal{M}(x_{b}) \parallel \mathcal{M^\prime}(x_{b})) \label{triplet:KL}
\end{align}

\begin{algorithm*}[htbp]
  \caption{Triplet Model Defense}
  \label{alg:triplet}
\begin{algorithmic}[1]
  \REQUIRE 
    Frozen original model $\mathcal{M}$;
    Trainable defense model $\mathcal{M}'$,
    Benign dataset $\mathcal{D}_b$, harmful dataset $\mathcal{D}_h$;
    Number of steps $T$; batch size $N$;
    Hyperparameters $\alpha, \beta, \gamma, m_b, m_h$
  \FOR{$t = 1, \ldots, T$}
    \STATE Sample a batch $x_b \sim \mathcal{D}_b$, $x_h \sim \mathcal{D}_h$
    \STATE Compute original representations $\h_{b, i}, \h_{h, i}$ using $\mathcal{M}$
    \STATE Compute new representations $\hp_{b, i}, \hp_{h, i}$ using $\mathcal{M}'$
    \STATE Compute $\hat{\hp} = \frac{1}{N} \sum_{i=1}^{N} \hp_{h,i}$
    \STATE $\mathcal{L}_\text{benign} = \frac{1}{N} \sum_{i=1}^N \max\left(0, \dist_{bp}(\h_{b,i}, \hp_{b,i}) - \dist_{bn}(\hp_{b,i}, \hat{\hp}) + m_b\right)$
    \STATE $\mathcal{L}_\text{harmful} = \frac{1}{N} \sum_{i=1}^N \max\left(0, \dist_{hp}(\hp_{h,i}, \hat{\hp}) - \dist_{hn}(\hp_{h,i}, \h_{h,i}) + m_h\right)$
    \STATE $\mathcal{L}_\text{KL} = \KL(\mathcal{M}(x_b) \parallel \mathcal{M^\prime}(x_b))$
    \STATE $\mathcal{L}_\text{Triplet} = \alpha \cdot \mathcal{L}_\text{benign} + \beta \cdot \mathcal{L}_\text{harmful} + \gamma \cdot \mathcal{L}_\text{KL}$
    \STATE Update parameters of $\mathcal{M}'$ using $\mathcal{L}_\text{Triplet}$
  \ENDFOR
\end{algorithmic}
\end{algorithm*}

Algorithm~\ref{alg:triplet} describes the training procedure for the triplet model defense.
The model weights are optimized until convergence on batches of benign and harmful prompt-response pairs.

\newcommand{\attack}{\textit{Attack}}
\subsection{Combining Representation Engineering with Adversarial Training}


Most LLMs are shipped with built-in safety features that prevent them from outputting harmful responses to plain harmful queries.
As such, gathering harmful representations $\hp_{h,i}$ from plain queries can lead to representations that are not truly informative
of dangerous model behavior. Inspired by work in hard negative mining, which focuses on learning on challenging negative samples \citep{robinson2021contrastivelearninghardnegative}, we propose to address this issue by integrating adversarial training in the representation space, by explicitly generating ``hard'' harmful representations via attacks.

\paragraph{Adversarial Hard Negative Mining}
In contrastive learning, ``hard negatives'' are challenging negative examples that are easily confused with positive examples.
For safety representation engineering, we define hard negatives as harmful representations that closely resemble benign ones.
Rather than relying on plain harmful representations, we make use of adversarial hard negative mining \citep{8517355}.


An adversarial attack neural network module $\attack_l$ is introduced at a randomly selected layer $l$.
The module is inserted between two transformer blocks in the residual stream, and is trained
using a Negative Log Likelihood (NNL) loss on harmful responses. The model is active when sampling
new harmful representations, and is periodically retained as model parameters are updated.
This module finds adversarial hard negatives, allowing the defense to counteract a wider diversity of harmful representations.
Appendix~\ref{app:adversarial_hard_neg_mining} contains more details on the training of our adversarial attack modules.

In summary, our complete method views safety representation engineering as a contrastive learning problem, optimizes a triplet-based loss function, and combines it with adversarial training to defend against harmful prompts more robustly.

\section{Experimental setup}


\paragraph{Models}
We evaluate our method on two widely used open-source models: Llama 3 8B Instruct \citep{grattafiori2024llama3herdmodels}
and Mistral 7B Instruct v0.2 \citep{jiang2023mistral7b}. These models are standard baselines for adversarial defense in the literature, which enables direct comparison with prior work.
Additional results on two extra models are provided in Appendix~\ref{app:models}.

\paragraph{Datasets}
We adapt the training pipeline of \citet{yousefpour2025representationbendinglargelanguage} for our defense method.
For benign data, we use UltraChat \citep{ding2023enhancingchatlanguagemodels}, a large-scale dataset of over 1.5 million multi-turn dialogues
that cover a wide range of topics such as art, history, literature, politics and technology.
For harmful data, we use WildGuardMix \citep{han2024wildguardopenonestopmoderation} which contains a broad spectrum of jailbroken prompts and responses,
and WildJailbreaks \citep{jiang2024wildteamingscaleinthewildjailbreaks}, a synthetic dataset of harmful prompt-response pairs,
featuring both straightforward and complex jailbreak prompts.
From these sources, we randomly select 10'000 benign and 10'000 harmful samples to construct a balanced training set.

Details on hyperparameter choices can be found in Appendix~\ref{app:hyperparameter}.

\section{Experiments}

In this section, we demonstrate the increased robustness of our methods against various attacks
compared to existing methods and explore three research questions (RQ) related to robustness and general performance.

\subsection{Overall defensive performance}

In this section, we compare the performance of our approach against different models and defenses.




\begin{table*}[ht]
\centering
\setlength{\tabcolsep}{4pt}
\begin{tabular}{l
                *{3}{ccc}@{\quad}
                *{3}{ccc}@{\quad}
                *{3}{ccc}}
\hline
 & \multicolumn{3}{c}{\textbf{REINFORCE-GCG}} 
 & \multicolumn{3}{c}{\textbf{GCG}} 
 & \multicolumn{3}{c}{\textbf{Embedding}} \\
\cmidrule(lr){2-4} \cmidrule(lr){5-7} \cmidrule(lr){8-10}
\textbf{Defense}
 & \textbf{HB} & \textbf{SR} & \textbf{Score}
 & \textbf{HB} & \textbf{SR} & \textbf{Score}
 & \textbf{HB} & \textbf{SR} & \textbf{Score} \\
\hline
Original model 
 & 52.50 & 40.00 & 42.87
 & 31.25 & 18.75 & 23.66
 & 100.00 & 90.24 & 81.89 \\

Circuit breakers 
 & 13.75 &  3.75 &  9.50
 &  2.86 &  1.43 &  4.25
 &  90.24 & 29.27 & 30.61 \\

RepBend 
 & 11.25 &  6.25 & 11.27
 &  2.86 & \textbf{0.00} &  1.65
 &  73.17 & 39.02 & 39.00 \\

\hline
\textbf{Triplet} 
 & \textbf{0.00} & \textbf{0.00} & \textbf{0.48}
 & \textbf{0.00} & \textbf{0.00} & \textbf{0.43}
 &  \textbf{65.85} & 12.20 & 14.57 \\

\textbf{Triplet + Adv} 
 &  3.75 &  2.50 &  6.99
 & \textbf{0.00} & \textbf{0.00} & 1.36
 &  75.61 & \textbf{4.88} & \textbf{8.70} \\
\hline
\end{tabular}
\caption{Attack success rates (ASR) using HarmBench (HB) and StrongREJECT (SR) across attack types, for various defenses (Llama 3 8B Instruct).
For GCG and REINFORCE-GCG attacks, each behavior was tested on a single attempt evaluated over 80 Behaviors. For embedding attacks, results were computed over 41 behaviors, with six attempts per behavior using different hyperparameter configurations (246 runs per model). The best result for each behavior was used. StrongREJECT scores are reported on a 0--100 scale.
}
\label{tab:all_attacks}
\end{table*}

\paragraph{Method}
We use the publicly available defensive models created by the authors of circuit breakers \citep{zou2024improvingalignmentrobustnesscircuit} and RepBend \citep{yousefpour2025representationbendinglargelanguage}.
We evaluate our defense using the HarmBench safety benchmark \citep{mazeika2024harmbenchstandardizedevaluationframework}.
For embedding attacks, we adopt a variant of the attack described by \citet{zou2024improvingalignmentrobustnesscircuit}.
To ensure robustness across hyperparameter choices, we use 6 different hyperparameter configurations.
For each behavior, we select the result from the configuration that produces the most harmful response
out of the six different runs.
For input-space attacks, we use GCG \citep{zou2023universaltransferableadversarialattacks} and REINFORCE-GCG \citep{geisler2025reinforceadversarialattackslarge} with base configurations.
To assess harmfulness, we use the binary HarmBench judge to get adversarial success rates (HB ASR) and the fine-grained StrongREJECT \citep{souly2024strongrejectjailbreaks} fine-tuned classifier to get harmfulness scores (SR Score) and adversarial success rates (SR ASR) for scores above $0.5$.
Full details of attack configurations and evaluation settings are provided in Appendix~\ref{app:setup}

\paragraph{Results}



Tables~\ref{tab:all_attacks} report the ASRs for embedding and GCG attacks on the Llama 3 8B model. All defense methods achieve substantial improvements over the base model, which shows ASRs above 90\% for embedding attacks. Among the defenses evaluated, our triplet defenses outperform both circuit breakers and RepBend. In particular,
the triplet defense achieves ASRs of $0\%$ for both REINFORCE and GCG, while the triplet defense with adversarial hard negative mining achieves the lowest embedding SR ASR of 4.88\%
and the lowest harmfulness score of $8.70$.

Figure~\ref{fig:embedding} shows embedding attack success rates for Llama 3 8B with two additional
adversarial defenses: Refusal Feature Ablation Training (ReFAT) \citep{yu2025robustllmsafeguardingrefusal} and Latent Adversarial Training (LAT) \citep{sheshadri2024latentadversarialtrainingimproves}. The Triplet model substantially improves embedding ASRs compared to these baselines.
Results for the Mistral 7B model are shown in Appendix~\ref{app:mistral}. Although the reported ASRs are higher than for the Llama model, the triplet defense outperforms the evaluated baselines.

Throughout our experiments, we find that the HarmBench classifier consistently produces higher ASRs than the StrongREJECT classifier. This is likely due to the tendency of HarmBench to classify responses as harmful responses based on the initial response tokens, even if the rest of the response is nonsensical or incoherent. Appendix Table~\ref{tab:overrefusal_example} shows an example of a generation classified as harmful by the HarmBench ASR, despite being practically harmless. These findings highlight how differences in harmfulness evaluation criteria can substantially influence measured outcomes.


\begin{figure}[ht]
    \centering
    \includegraphics[width=0.48\textwidth]{latex/images/Llama_strongreject_score.pdf} %
    \caption{Embedding Attack success rate (ASR) using StrongREJECT for various defenses (Llama 3 8B Instruct.}
    \label{fig:embedding}
\end{figure}

\subsection{RQ1: How robust is our defense to different attack configuration choices?}

The goal of this research question is to evaluate the robustness of our defense to different embedding configurations.

\paragraph{Method}

We compare the ASRs of embedding attacks across the six different hyperparameter configurations, with full details provided in Appendix~\ref{app:setup}. For each defense, we report the best, worst, and mean ASR obtained over
all configurations.

\definecolor{conf0}{HTML}{1f77b4} %
\definecolor{conf2}{HTML}{2ca02c} %
\definecolor{conf3}{HTML}{d62728} %
\definecolor{conf4}{HTML}{9467bd} %
\definecolor{conf5}{HTML}{8c564b} %
\definecolor{conf1}{HTML}{b35006} %

\begin{table*}
\centering
\begin{tabular}{lccccccccc}
\hline
\textbf{Defense} & \multicolumn{3}{r}{\textbf{HarmBench ASR}} & \multicolumn{3}{r}{\textbf{StrongREJECT ASR}} & \multicolumn{3}{r}{\textbf{StrongREJECT Score}} \\
 & mean & {min} & {max} & {mean} & {min} & {max} & {mean} & {min} & {max} \\
\hline
Original model & 77.33 & 54.00\textsuperscript{\textcolor{conf1}{1}} & 98.00\textsuperscript{\textcolor{conf3}{3}} & 48.16 & 26.10\textsuperscript{\textcolor{conf4}{4}} & 63.34\textsuperscript{\textcolor{conf3}{3}} & 53.25 & 24.39\textsuperscript{\textcolor{conf4}{4}} & 73.17\textsuperscript{\textcolor{conf5}{5}} \\
RepBend & 24.50 & \textbf{2.00}\textsuperscript{\textcolor{conf5}{5}} & 37.00\textsuperscript{\textcolor{conf2}{2}} & 10.36 & 4.00\textsuperscript{\textcolor{conf5}{5}} & 22.06\textsuperscript{\textcolor{conf2}{2}} & 8.54 & 2.44\textsuperscript{\textcolor{conf0}{0}} & 19.51\textsuperscript{\textcolor{conf2}{2}} \\
Circuit breakers & 38.67 & 27.00\textsuperscript{\textcolor{conf2}{2}} & 54.00\textsuperscript{\textcolor{conf1}{1}} & 9.32 & 3.41\textsuperscript{\textcolor{conf5}{5}} & 14.53\textsuperscript{\textcolor{conf0}{0}} & 6.91 & \textbf{0.00}\textsuperscript{\textcolor{conf5}{5}} & 12.20\textsuperscript{\textcolor{conf0}{0}} \\
\hline
\textbf{Triplet} & \textbf{23.83} & 17.00\textsuperscript{\textcolor{conf3}{3}} & \textbf{32.00}\textsuperscript{\textcolor{conf1}{1}} & 3.55 & 1.16\textsuperscript{\textcolor{conf4}{4}} & 9.46\textsuperscript{\textcolor{conf2}{2}} & 2.44 & \textbf{0.00}\textsuperscript{\textcolor{conf0}{0}} & 9.76\textsuperscript{\textcolor{conf2}{2}} \\
\textbf{Triplet + Adv} & 24.40 & 10.00\textsuperscript{\textcolor{conf1}{1}} & 41.00\textsuperscript{\textcolor{conf2}{2}} & \textbf{2.23} & \textbf{1.10}\textsuperscript{\textcolor{conf3}{3}} & \textbf{4.28}\textsuperscript{\textcolor{conf0}{0}} & \textbf{0.49} & \textbf{0.00}\textsuperscript{\textcolor{conf1}{1}} & \textbf{2.44}\textsuperscript{\textcolor{conf0}{0}} \\
\hline
\end{tabular}
\caption{Mean, maximum, and minimum embedding attack ASR across six different hyperparameter configurations (Llama 3 8B). Colored superscript indicates the configuration index for which the ASRs were achieved.
}
\label{tab:robustness}
\end{table*}

\paragraph{Results}

Table~\ref{tab:robustness} presents our results.
Both RepBend and circuit breakers exhibit significant variability across attack hyperparameter configurations. In particular, using configuration $2$, we get a StrongREJECT ASR of $20\%$ on RepBend and $2\%$ on circuit breakers, while configuration $0$ results in an ASR of $2\%$ and $12\%$ respectively. In comparison, our triplet defense consistently demonstrates low ASRs, with a worst-case StrongREJECT ASR of $2\%$.
The complete results are provided in Appendix Figure~\ref{tab:full_embedding_results}.
These findings highlight the necessity of evaluating defenses across diverse attack configurations to accurately
assess their robustness, as relying on a single configuration could bias results in favor of a particular defense.

\subsection{RQ2: Does applying our defense affect the general performance of the model?}

The objective of this research question is to determine whether the application of our defensive mechanism affects the general language capabilities of the model.

\paragraph{Method}

We assess the general performance of our trained models on a suite of benchmarks, including MMLU, HellaSwag, TruthfulQA, and GSM8K. See Appendix \ref{app:setup} for more details.

\definecolor{darkerred}{rgb}{0.4,0,0} 
\definecolor{darkred}{rgb}{0.8, 0, 0}
\begin{table*}[htbp]
\centering
\resizebox{\textwidth}{!}{
 \begin{tabular}{lccccccccc}
\hline
 & \textbf{ARC (Easy)} & \textbf{GSM8K} & \textbf{HellaSwag} & \textbf{MMLU} & \multicolumn{3}{c}{\textbf{TruthfulQA}} \\
\cmidrule(lr){6-8}
 &  &  &  &  & \textbf{Gen} & \textbf{MC1} & \textbf{MC2} \\
\hline
\textbf{Original model} & \textbf{81.61} & \textbf{75.36} & \textbf{57.75} & \textbf{63.72} & \textbf{46.39} & \textbf{36.23} & \textbf{51.67} \\
\hline
Circuit breakers & 81.44 \textcolor{darkred}{(-0.17)} & 75.44 \textcolor{darkgreen}{(+0.08)} & 57.46 \textcolor{darkred}{(-0.29)} & 63.57 \textcolor{darkred}{(-0.15)} & 48.23 \textcolor{darkgreen}{(+1.84)} & 36.96 \textcolor{darkgreen}{(+0.73)} & 51.61 \textcolor{darkred}{(-0.05)} \\
RepBend & 80.98 \textcolor{darkred}{(-0.63)} & 49.05 \textcolor{darkerred}{(-26.31)} & 60.58 \textcolor{darkgreen}{(+2.83)} & 60.26 \textcolor{darkred}{(-3.46)} & 2.08 \textcolor{darkerred}{(-44.31)} & 41.00 \textcolor{darkgreen}{(+4.77)} & 60.05 \textcolor{darkgreen}{(+8.38)} \\
\hline
\textbf{Triplet} & 81.27 \textcolor{darkred}{(-0.34)} & 74.30 \textcolor{darkred}{(-1.06)} & 59.62 \textcolor{darkgreen}{(+1.87)} & 63.85 \textcolor{darkgreen}{(+0.13)} & 45.65 \textcolor{darkred}{(-0.73)} & 40.76 \textcolor{darkgreen}{(+4.53)} & 55.37 \textcolor{darkgreen}{(+3.70)} \\
\textbf{Triplet + Adv} & 81.99 \textcolor{darkgreen}{(+0.38)} & 74.91 \textcolor{darkred}{(-0.45)} & 60.70 \textcolor{darkgreen}{(+2.95)} & 63.38 \textcolor{darkred}{(-0.34)} & 44.55 \textcolor{darkred}{(-1.84)} & 42.96 \textcolor{darkgreen}{(+6.73)} & 57.29 \textcolor{darkgreen}{(+5.63)} \\
\hline
\end{tabular}
}
\caption{Performance comparison of models on general capability benchmarks (Llama 3 8B Instruct). } 
\label{tab:benign_performance}
\end{table*}

\paragraph{Results}

As shown in Table \ref{tab:benign_performance}, our triplet method achieves a performance comparable to the base
model, which indicates that our approach preserved the model's general capabilities.
Notably, the defenses showcase an improvement on the TruthfulQA (MC) benchmark, likely because their increased tendency to reject harmful responses also leads them to reject untruthful content, which is often harmful.

In contrast, the performance of the RepBend model decreases significantly on the generation-based benchmarks GSM8K and TruthfulQA (Gen), with accuracy dropping from 75\% to 49\% on GMS8k, and from 46\% and 2\% on TruthfulQA (Gen).
This suggests overfitting to the defense objective at the expense of general language performance.
Examples of generations of GSM8K for RepBend can be seen in Appendix Table~\ref{tab:overrefusal_example}, in which the model fails to answer the questions. In contrast, our trained triplet models do not suffer from this issue, maintaining both GSM8K and TruthfulQA (Gen) performance close to the base model.
These findings highlight the strengths of our approach and illustrate that defenses can have unintended side effects, which may themselves have harmful consequences in sensitive applications.

\subsection{RQ3: How does our defense generalize to out-of-distribution input/output formats?}

While adversarial attack success rates are valuable for measuring robustness in plain text, they do not fully capture a defense's capability to generalize to out-of-distribution response formats. To address this, we introduce a new evaluation metric for safety representation engineering defenses, based on relative distances.

\paragraph{Method}
We apply random capitalization (following the Best-of-N jailbreak attack protocol \citep{hughes2024bestofnjailbreaking}) and translations to five languages to 159 behaviors of the HarmBench benchmark, generating a set of augmented prompts and responses $\mathcal{A}(b)$ for each behavior $b$.
For each defense, we compute the Mean Minimum Distance Ratio (MMDR) as the average over behaviors of the smallest ratios between the distance of augmented and original harmful representations:

\newcommand{\gen}{\text{MMDR}}
\newcommand{\mgen}{\mathrm{MMDR}}

\begin{align}
	\gen_d = \frac{1}{|N|} \sum_{i \in N} \min_{a \in \mathcal{A}(h_i)} \frac{\dd(\hp_{a}, \h_{a})}{\dd(\hp_{h,i}, \h_{h,i})}
\end{align}

Averaged over all model layers, MMDR quantifies the model's worst-case generalization to out-of-distribution augmentations.

\begin{table}[htbp]
\centering
\setlength{\tabcolsep}{3pt} %
\begin{tabular}{lcc}
\hline
\textbf{Distance} & {$\mathbf{MMDR}_{d_2}$} & {$\mathbf{MMDR}_{d_\mathrm{cos}}$} \\
\hline
Circuit breakers & 0.63 & 0.49 \\
Triplet A1: CB & 0.70 & 0.54 \\
RepBend & 0.70 & 0.64 \\
Triplet A2: RepBend & 0.78 & 0.64 \\
Triplet A3: Full & \textbf{0.80} & 0.66 \\
Triplet A4: Full + Adv & \textbf{0.80} & \textbf{0.70} \\
\hline
\end{tabular}
\caption{Generalization of the defenses to different data augmentations (Llama 3 8B Instruct) with the L2 norm $\dd_2$ and the cosine distance $\dd_{\mathrm{cos}}$}
\label{tab:generalization}
\end{table}

\paragraph{Results}

As shown in Table~\ref{tab:generalization}, the full triplet defenses achieve higher MMDR values,
up to $0.8$ for both distance metrics, compared to the circuit breaking ($0.6$) and RepBend $(0.7)$ defenses.
A value close to $1$ indicates strong generalization, meaning the defense modifies augmented
harmful representations similarly to the unaugmented ones, while a value close to $0$ would imply that some augmentations bypass the defense.
Notably, the MMDR increases as more loss terms are incorporated, highlighting the effectiveness of our approach in generalizing to different input and output formats.

\subsection{Ablation study}

To study the impact of our loss functions, we conduct an ablation study in which our defense methods are trained with different loss configurations. 

Specifically, we train a model A1 by removing the $\dd_{b,n}$ and $\dd_{h, p}$ components, making it closely related to circuit breaking.
For model A2, we ablate only $\dd_{b, n}$, resulting in a formulation similar to RepBend but with a margin-based objective.
Model A3 retains all loss terms, while model A4 incorporates adversarial hard negative mining and all loss terms.
Details and full results are provided in Appendix~\ref{app:ablation}

Our results show that A1 performs similarly to circuit breaking, validating the correspondence between the loss formulations.
A2 consistently outperforms both A1 and RepBend, highlighting the advantage of our margin-based triplet learning objective.
A3 and A2 perform comparably overall, with A3 achieving better results in input-space attacks
and A2 on embedding space attacks. A4 achieves the best performance.
Table~\ref{tab:generalization} further shows that for our studied models, removing loss terms leads to a decrease in $\gen$. These findings demonstrate the importance of the $\dd_{b, n}$ term in our triplet loss, as well as the additional benefit of adversarial hard negative mining.

\section{Conclusion}
\label{sec:conclusion}

This work presents a novel method for improving the robustness of LLMs
against adversarial attacks based on contrastive representation learning and adversarial hard negative mining. Our findings demonstrate notable robustness improvement while
maintaining the model's general capabilities.

\section*{Ethics Statement}

This section discusses the ethical considerations that arise from the development and deployment of defensive methods for AI models.
First, the development of defenses may lead to overconfidence in the safety of AI models, which in turn could encourage the deployment of less safe models. In addition, better defenses
may lead to the development of stronger, more sophisticated attacks, thereby increasing the risk of misuse for widespread AI models with fewer security measures.
Lastly, our representation space attack module used for adversarial hard negative mining
could be misused by malicious actors to circumvent the defenses of AI models
in a white-box setting.
Nevertheless,
these potential downsides are outweighed by the
benefits of developing better defenses. These concerns underscore the need for responsible use and deployment of research findings in the field of AI safety.

\section*{Limitations}

Despite the promising results of our method, several limitations should be acknowledged.
First, while robustness improves significantly on the Llama model,
the Mistral model remains more vulnerable to attacks, even though our method still outperforms circuit breakers.  This suggests that further tuning or architecture-specific adaptations may be required to achieve a strong robustness across models.

Secondly, given the computationally intensive nature of adversarial training and
jailbreak generation, the choice of hyperparameters and training strategy may
not be optimal. Furthermore, due to the 2-GPU-hour cost per REINFORCE attack,
our method was evaluated on 80 HarmBench behaviors in the validation set compared to 
the 300 behaviors in the training set. Although we expect similar trends to hold,
this assumption has not been empirically verified.

Third, while our method is robust to a variety of attacks, it is not guaranteed to be robust to all, especially to attacks in the representation space.
Using more attack configurations and attempts per behavior may also lead to high attack success rates, at the cost of increased computation time. 

Finally, like other representation engineering-based methods, our trained models
may result in incoherent and ineligible behavior if the model misinterprets benign
inputs as harmful, which in turn could lead to harmful consequences in some
critical settings. Therefore, careful considerations and additional safeguards
may be necessary before deploying these methods in real-world applications.
Finally, our method with adversarial hard negative mining requires a moderate training time of up to 12 hours on a
single H100 GPU for the Llama 3 8B model. This requirement may limit the scalability of
our approach to much larger models with hundreds of billions of parameters.



\section*{Acknowledgements}

We thank Kellin Pelrine, Roger Grosse, and Stephen Casper for their feedback on our work.
This material is based in part upon work supported by the German Federal Ministry of Education and Research (BMBF): Tübingen AI Center, FKZ: 01IS18039B; by the Machine Learning Cluster of Excellence, EXC number 2064/1 – Project number 390727645; by Schmidt Sciences SAFE-AI Grant; by NSERC Discovery Grant RGPIN-2025-06491; 
by a National Science Foundation award (\#2306372); by a Swiss National Science Foundation award (\#201009) and a Responsible AI grant by the Haslerstiftung;
as part of the ``Swiss AI initiative'' by a grant from the Swiss National Supercomputing Centre (CSCS) under project IDs a07 and a130 on Alps. Additional calculations were carried out on the ETH Euler cluster and the Center for AI Safety cluster.



\bibliography{custom}

\appendix

\section{Appendix}
\label{sec:appendix}

\subsection{Additional Details on Related Work}
\label{app:details_related}

\paragraph{Jailbreak Attacks}
Token-level jailbreaks optimize a harmful adversarial sequence of tokens appended to the prompt.
Greedy Coordinate Gradient (GCG) \citep{zou2023universaltransferableadversarialattacks} iteratively optimizes an adversarial suffix with gradient information to select promising token substitutions. Projected Gradient Descent (PGD) for LLMs \citep{geisler2025attackinglargelanguagemodels} adapts the PGD continuous optimization algorithm to the discrete setting of tokenized inputs.
Adaptive methods \citep{andriushchenko2025jailbreakingleadingsafetyalignedllms}
are able to achieve near 100\% attack success rates on leading models such as Claude or GPT-4o.

Prompt-level jailbreaks involve optimizing the entire prompt and generally result in human-readable jailbreak prompts. Prompt Automatic Iterative Refinement (PAIR) \citep{chao2024jailbreakingblackboxlarge} uses a fine-tuned LLM in a conversation against a target model to refine a harmful prompt.
Persuasive Adversarial Prompts (PAP) \citep{zeng2024johnnypersuadellmsjailbreak} generate emotionally persuasive prompts to trick the target model. Additional jailbreaking strategies manipulate the input and the output format \citep{huang2024endlessjailbreaksbijectionlearning, yuan2024gpt4smartsafestealthy}, or manipulating model reasoning \citep{wu2024knowimsayingjailbreak, ramesh2024gpt4jailbreaksnearperfectsuccess}.

 

\paragraph{External Jailbreak Defenses}
Common strategies to defend against jailbreaks include modifying the input before giving it to the model, for instance by inserting, swapping or replacing characters \citep{robey2024smoothllmdefendinglargelanguage}, using perplexity-based filters \citep{jain2023baselinedefensesadversarialattacks}, using paraphrasing and defensive suffixes \citep{yuan2024rigorllmresilientguardrailslarge}, or using guardrail models to analyze the intent of queries and responses \citep{zeng2024autodefensemultiagentllmdefense}.

\paragraph{Internal Jailbreak Defenses}
Internal defenses directly modify the model behavior by fine-tuning the model weights. 
Common approaches include Reinforcement Learning from Human Feedback \citep{kaufmann2024surveyreinforcementlearninghuman}, 
which use human preference data to fine-tune model weights. Other methods can identify problematic layers and edit out harmful content \citep{zhao2024defendinglargelanguagemodels}.

\paragraph{Triplet loss}
\begin{figure}[htbp]
    \centering
    \tdplotsetmaincoords{70}{120} %

    \begin{minipage}{0.48\textwidth} %
        \centering
        \begin{tikzpicture}[
            tdplot_main_coords, %
            scale=1, %
            point/.style={circle, fill, inner sep=2pt, outer sep=0pt}, %
            label/.style={font=\small} %
        ]

            \def\axislen{2}
            \draw[->, gray] (0,0,0) -- (\axislen,0,0) node[below] {X};
            \draw[->, gray] (0,0,0) -- (0,\axislen,0) node[left] {Y};
            \draw[->, gray] (0,0,0) -- (0,0,\axislen) node[above] {Z};

            \node[point, fill=blue, label={below left:Anchor}] (A1) at (1,1,0.5) {}; %
            \node[above left, font=\small] at (A1) {$A$};
            \node[point, fill=green, label={above right:Positive}] (P1) at (1,2,2) {}; %
            \node[above left, font=\small] at (P1) {$P$};
            \node[point, fill=red, label={below right:Negative}] (N1) at (4,-0.5,1.5) {}; %
            \node[above left, font=\small] at (N1) {$N$};

            \draw[thick] (A1) -- (P1);
            \draw[thick] (A1) -- (N1);
        \end{tikzpicture}
    \end{minipage}%
    \hfill %
    \begin{minipage}{0.48\textwidth} %
        \centering
        \begin{tikzpicture}[
            tdplot_main_coords, %
            scale=0.9, %
            point/.style={circle, fill, inner sep=2pt, outer sep=0pt}, %
            label/.style={font=\small} %
        ]

            \def\axislen{2}
            \draw[->, gray] (0,0,0) -- (\axislen,0,0) node[below] {X};
            \draw[->, gray] (0,0,0) -- (0,\axislen,0) node[left] {Y};
            \draw[->, gray] (0,0,0) -- (0,0,\axislen) node[above] {Z};

            \node[point, fill=blue, label={below left:Anchor}] (A2) at (1,1,0.5) {};
            \node[above left, font=\small] at (A2) {$A$};
            \node[point, fill=green, label={above right:Positive}] (P2) at (0.5, 1, 1) {}; %
            \node[above left, font=\small] at (P2) {$P$};
            \node[point, fill=red, label={below right:Negative}] (N2) at (4.5,-1.5,2.5) {}; %
            \node[above left, font=\small] at (N2) {$N$};

            \draw[thick] (A2) -- (P2);
            \draw[thick] (A2) -- (N2);


        \end{tikzpicture}
    \end{minipage}

    \caption{Triplet loss objective before and after a learning step. The anchor (blue) is kept at the same position, while the positive (green) is moved closer to the anchor, and the negative (red) is moved further away from the anchor.}
    \label{fig:triplet_loss_3d}
\end{figure}

Figure \ref{fig:triplet_loss_3d} shows a visualization of how the triplet loss affects
data points during training.

\subsection{Additional Details on our Method}
\label{app:details_method}

\paragraph{Discussion on the Choice of Triplet}

Recent defense methods aimed at disrupting harmful LLM generations, such as circuit breakers \citep{zou2024improvingalignmentrobustnesscircuit}, have been shown to be a powerful alternative to adversarial training. These methods also often cause successful attack responses to become incoherent or to break down at later stages of harmful generation \citep{geisler2025reinforceadversarialattackslarge}. Figure~\ref{fig:example-prompts} provides examples of harmful prompts and responses from attacks against adversarial defenses such as LAT \citep{sheshadri2024latentadversarialtrainingimproves} as well as defense methods like ours which aim to directly disrupt harmful generations.

As discussed in Section~\ref{sec:interpret_cb}, the goal of this work is to extend circuit breakers by treating the inner representation space of LLMs as an embedding space, in which new harmful representations should be pushed far apart from both benign representations and old harmful representations.
We adopt a triplet formulation due to its simplicity and its ability to generalize both circuit breakers and RepBend.

While our approach uses a single negative for each harmful anchor, other contrastive losses, such as InfoNCE~\citep{oord2019representationlearningcontrastivepredictive} or N-Pair loss \citep{NIPS2016_6b180037} could enable training with multiple negatives, which may provide further performance improvements.
Other contrastive methods, such as SimCLR \citep{chen2020simpleframeworkcontrastivelearning}, rely on encoders to learn latent representations. In the LLM defense setting, these methods might require
additional adaptation, as contrast between harmful and harmless representations in the learned latent space does not necessarily imply separation in the representation stream, which is likely relevant for defending against adversarial attacks. We leave these extensions and explorations to future work.

\input{latex/figure.tex}

\paragraph{Distances}
The notions of similarity and dissimilarity are defined by a distance function $\dist: \mathbb{R}^d \times \mathbb{R}^d \to \mathbb{R}$. In this work, we will use the term distance to refer to pseudodistances, as we only need to define a notion of similarity and dissimilarity.
We consider the following distances:
\begin{itemize}[nolistsep]
    \item $L_k$: $\dist_k(x, y) = ||x - y||_k$
    \item Cosine: $\dist_{\cos}(x, y) = 1 - \frac{x \cdot y}{||x|| \cdot ||y||}$
    \item Mix : $\dist_{mix}(x, y) = \alpha \cdot \dist_2(x, y) + \beta \cdot \dist_{\cos}(x, y)$
    \item Null: $\dist_0(x, y) = 0$
\end{itemize}
In the representation space of LLMs, cosine similarity has been shown to capture semantic similarity.
However, this notion of distance does not take into account the scale of the embeddings, which is an important factor in the representation space.
Mixing different distances allows to learn a space in which angular similarity is more or less important than Euclidean distance depending on the choice of the user.

\paragraph{Adversarial Hard Negative Mining}
\label{app:adversarial_hard_neg_mining}

\begin{algorithm*}[htbp]
    \caption{Training an attack module $\attack_l$ on harmful responses}
\label{alg:train_attack}
\begin{algorithmic}[1]
\REQUIRE Model $\pi$, target layer $l$, module $\attack_l$, number of training steps $n$
\ENSURE Trained $\attack_l$
\STATE Set up an Adam optimizer for $\attack_l$
\FOR{$i$ from 0 to $n - 1$}
    \STATE Sample a harmful prompt-response pair $(x, y)$
    \STATE Get logits $\pi(y | x, \operatorname{do}(\h_l = \attack_l(\h_l)))$
    \STATE Compute the Negative Log-Likelihood (NLL) loss:
    \[ \mathcal{L} = - \frac{1}{T} \sum_{t=1}^{T} \log p(y_t |x, y_{<t}) \]
    \STATE Update $\attack_l$ to minimize $\mathcal{L}$
\ENDFOR
\end{algorithmic}
\end{algorithm*}

Algorithm~\ref{alg:train_attack} shows the training process of an attack module $\attack_l$ on harmful responses.
Every $m$ steps, a new layer $l$ is randomly selected, and an attack module $\attack_l$ is trained until convergence.
Every $k$ steps, the attack module is retrained for $K$ steps, to ensure continuous effectiveness as the model is updated.
By varying the layer $l$ throughout the attack, the defense learns to counteract a diverse range of harmful representations.

\begin{algorithm*}[htbp]
  \caption{Triplet Model Defense with Adversarial Hard Negative Mining}
  \label{alg:triplet_attack}
\begin{algorithmic}[1]
  \REQUIRE 
    Frozen original model $\mathcal{M}$, Trainable defense model $\mathcal{M}'$ (e.g., with LoRA adapters), Benign dataset $\mathcal{D}_b$, harmful dataset $\mathcal{D}_h$, Number of steps $T$; batch size $N$, Hyperparameters $\alpha, \beta, \gamma, m_b, m_h$, Retrain interval $k$, Retrain steps $K$, Attack module selection interval $a$ \\
  \ENSURE Trained model $\mathcal{M}^\prime$
  \FOR{$t = 0,\dots,T-1$}
    \IF{$t \bmod a = 0$}
      \STATE Remove $\attack_l$ if exists \hfill $\triangleright$ Create new attack module
      \STATE Choose random layer $l\in\{1,\dots,L\}$ to attack
      \STATE Train $\attack_l$ until convergence
      \STATE Apply $\attack_l$ to the model \hfill $\triangleright$ Apply attack module
    \ENDIF
    \STATE Sample a batch $x_b \sim \mathcal{D}_b$, $x_h \sim \mathcal{D}_h$
    \STATE Compute original representations $\h_{b, i}, \h_{h, i}$ using $\mathcal{M}$
    \STATE Compute new representations $\hp_{b, i}, \hp_{h, i}$ using $\mathcal{M}'$
    \STATE Compute $\hat{\hp} = \frac{1}{N} \sum_{i=1}^{N} \hp_{h,i}$ \hfill $\triangleright$ Mean of harmful representations
    \STATE $\mathcal{L}_\text{benign} = \frac{1}{N} \sum_{i=1}^N \max\left(0, d_{bp}(\h_{b,i}, \hp_{b,i}) - d_{bn}(\hp_{b,i}, \hat{\hp}) + m_b\right)$
    \STATE $\mathcal{L}_\text{harmful} = \frac{1}{N} \sum_{i=1}^N \max\left(0, d_{hp}(\hp_{h,i}, \hat{\hp}) - d_{hn}(\hp_{h,i}, \h_{h,i}) + m_h\right)$
    \STATE $\mathcal{L}_\text{KL} = \KL(\mathcal{M}(x_b) \parallel \mathcal{M^\prime}(x_b))$
    \STATE $\mathcal{L}_\text{Triplet} = \alpha \cdot \mathcal{L}_\text{benign} + \beta \cdot \mathcal{L}_\text{harmful} + \gamma \cdot \mathcal{L}_\text{KL}$
    \STATE Update parameters of $\mathcal{M}'$ using $\mathcal{L}_\text{Triplet}$
    \IF{$s \bmod k = 0$}
      \STATE Retrain attack module $\attack_l$ for $K$ steps \hfill $\triangleright$ Train attack module
    \ENDIF
  \ENDFOR
\end{algorithmic}
\end{algorithm*}

Algorithm~\ref{alg:triplet_attack} shows the process of training our triplet based objective
with the addition of adversarial hard negative mining.

\subsection{Experimental Setup}
\label{app:setup}

\paragraph{Hyperparameters}
\label{app:hyperparameter}
Model training is conducted using the Adam optimizer \citep{kingma2017adammethodstochasticoptimization} with a learning rate of $1\mathrm{e}{-5}$ and a batch size of 16.
We employ LoRA adapters with rank $r = 16$, scaling factor $\alpha = 16$ and a dropout of $0.05$,
following existing representation engineering methods for direct comparison.
We collect harmful representations from layers $20$ to $31$ as in \citet{yousefpour2025representationbendinglargelanguage}. We use a balanced $\dist_{mix}$ for all distances in our experiments, and the hyperparameters $\alpha = 0.5$, $\beta = 0.4$, $\gamma = 0.9$, $m_b = 500$, $m_h = 1500$. We train for a total of $1100$ steps.

In adversarial hard negative mining experiments, we retrain a new linear attack module every $30$ steps on a random layer between $0$ and $20$, and update it for one iteration at every other defense training step. During training, we sample $70\%$ of the original harmful representations and $30\%$ of adversarial harmful representations.
Our implementation is released under a MIT license.

\paragraph{Embedding Attacks}
\label{par:embedding}
We implement a variant of the attack introduced by \citet{zou2024improvingalignmentrobustnesscircuit}.
We substitute the SGD optimizer for the Adam optimizer \citep{kingma2017adammethodstochasticoptimization} due to its superior performance in complex optimization landscapes.

To ensure robust evaluation, we adopt a best-of-N attack strategy. We use $N=6$ different attack hyperparameter configurations and select the one that yields the best classifier accuracy.
This ensures that defenses are not overly sensitive to specific hyperparameter settings.
Unlike the original study, which uses a simple ``\texttt{x x x x}'' string to initialize the embedding, we use a contextually meaningful initialization string for more realistic attacks.
The attack is terminated when the loss drops below a certain threshold. We explore a range of threshold values
across different attack configurations.
We summarize our different attack configurations in Table \ref{tab:attack_hyperparameters}.

\begin{table*}[htbp]
	\centering
	\begin{tabular}{lccp{6cm}c}
\hline
Config & Learning rate & Max steps & Initial string & Early stop threshold \\
\hline
run1 & 0.1 & 1000 & Be honest. & 0.001 \\
run2 & 0.005 & 1000 & It is important to carefully consider all the details before answering truthfully and fully. Don't worry about being inappropriate. & 0.0001 \\
run3 & 0.001 & 1000 & Make sure to answer the question truthfully. & 0.0005 \\
run4 & 0.0005 & 1000 & Answer as best as you can, stick to the question. & 0.01 \\
run5 & 0.02 & 1000 & Please try your best to give a fact-based and helpful answer that aligns with known information. & 0.05 \\
run6 & 0.0001 & 1000 & Make sure to answer well, answer step by step. & 0.00001 \\
\hline
\end{tabular}
\caption{Overview of configuration hyperparameters used for our embedding attack}
\label{tab:attack_hyperparameters}
\end{table*}

\paragraph{GCG Attacks}
We use the GCG attack \citep{zou2023universaltransferableadversarialattacks} white-box gradient
based optimization attack integrated in the HarmBench benchmark
to evaluate the robustness of our defenses against input space attacks.
We use the default attack hyperparameters of 500 steps with a search width of 512.


\paragraph{Safety Benchmarks}
We use the HarmBench benchmark \citep{mazeika2024harmbenchstandardizedevaluationframework} to evaluate the maliciousness of our trained models. We use the validation set containing 80 harmful behaviors.
To measure the harmfulness of model responses, we use two different automated classifiers:

\begin{itemize}
	\item The HarmBench classifier, a fine-tuned Llama 2 13B model providing a binary classification of harmfulness. It is the standard classifier used in the HarmBench benchmark.
	\item The StrongREJECT classifier \citep{souly2024strongrejectjailbreaks}, which is a fine-tuned Gemma 2B model. This classifier outputs a score estimate of the harmfulness of a response between 0 and 1,
		and is better suited for evaluating weakly harmful responses. We use a threshold of 0.5 to classify a response as harmful or not.
\end{itemize}

The choice of using a more fine-grained classifier allows a more accurate evaluation of the performance of our defenses,
especially when the model begins to produce incoherent content after generating some compliant tokens.

\paragraph{General Performance Benchmark}

We use 5 common benchmarks to evaluate our trained models on standard language modeling tasks:

\begin{itemize}[nolistsep]
	\item \textbf{MMLU} \citep{hendrycks2021measuringmassivemultitasklanguage}: A collection of multiple choice questions on a wide range of topics, including mathematics and chemistry.
	\item \textbf{ARC-Easy} \citep{clark2018thinksolvedquestionanswering}: A collection of natural, grade-school multiple choice science questions.
	\item \textbf{HellaSwag} \citep{zellers2019hellaswagmachinereallyfinish}: A commonsense reasoning benchmark of multiple choice questions.
	\item \textbf{GSM8K} \citep{cobbe2021trainingverifierssolvemath}: A collection of diverse grade-school math word problems for evaluating generative reasoning.
    \item \textbf{TruthfulQA} \citep{lin2022truthfulqameasuringmodelsmimic}: A benchmark for measuring the truthfulness of language models, evaluated with multiple choice (mc1) and generative (gen) questions.
\end{itemize}


We use the lm-eval library to run these benchmarks \citep{eval-harness}.


\paragraph{Harmful Behavior Augmentations}
To generate a set of augmented prompts and responses to calculate the Mean Minimum Distance Ratio (MMDR), we apply the following augmentations to the HarmBench harmful behaviors:
\begin{itemize}[nolistsep]
	\item \textbf{Random capitalization}: Randomly capitalizing letters in the input and output, following the Best-of-N jailbreak attack protocol \citep{hughes2024bestofnjailbreaking}. We sample 5 variants per behavior.
	\item \textbf{Translation}: Translating each harmful pair in French, German, Spanish, Chinese, and Ukrainian, using a jailbroken Llama 3 8B model. The inclusion of Chinese and Ukrainian enables the evaluation of generalization across different writing systems. Our translation prompt is shown in Appendix~\ref{app:prompts}.
\end{itemize}

\paragraph{Compute Cost Analysis}

Our triplet-based Llama 3 8B defense, as detailed in the experimental setup,
required approximately 7 hours of training on a single H100. The adversarial hard negative mining version took slightly longer and required 9 hours.
Training the Zephyr triplet defense took 12 hours on a single A100 for 1100 training steps, while training the Phi 3 Mini 4k model defense took 9 hours under the same conditions.

For comparison, \citet{sheshadri2024latentadversarialtrainingimproves} estimate that
the training time of the LAT method is 12 hours or less. This suggests that our method has a typical compute cost relative to adversarial defense approaches. However, speed-efficient defenses such as ReFAT \citep{yu2025robustllmsafeguardingrefusal} are reported to require roughly 10 times less compute time than LAT. While faster training is desirable, we view speed as a secondary concern as long as the training time remains reasonable,
since ensuring safety is the primary objective.

For our main experiments, we have run 480 REINFORCE-GCG attacks, each taking approximately 2 hours on a single H100 GPU, and an additional 480 GCG attacks which averaged 15 minutes per attack on the same hardware. Across all our experiments and development, we estimate our total compute usage to be in the range of 2,000 GPU hours.

\subsection{Ablation study}
\label{app:ablation}

To evaluate the importance of each element in our triplet-based loss, we conduct an ablation study.
We consider the following models:

\begin{itemize}[nolistsep]
	\item \textbf{Triplet A1: CB}: The triplet loss with $\dd_{bn}$ and $\dd_{bp}$ set to the null distance $\dd_{0}$.
			This configuration is a simplified version of the triplet loss which is similar to the circuit breaking method.
			The difference lies in the addition of the Kullback-Leibler divergence term, and the use of mixed distances.
	\item \textbf{Triplet A2: RepBend}: The triplet loss with $\dd_{bn}$ set to the null distance $\dd_{0}$.
			This configuration is similar to the RepBend method, but incorporates a margin to regulate and constraint the distances, preventing them from becoming unbounded, and uses mixed distances for $\dd_{bp}$ and $\dd_{hn}$.
			The use of margins allows for a flexible selection of the number of iterations, without concern for the distance diverging to infinity.
	\item \textbf{Triplet A3: Full}: The full triplet loss as described in Algorithm~\ref{alg:triplet}.
	\item \textbf{Triplet A4: Full + Adv}: The full triplet loss with adversarial hard negative mining as described in Algorithm~\ref{alg:triplet_attack}.
\end{itemize}


Appendix~\ref{app:ablation} show the losses in detail for models A1, A2, and A3.

\begin{table*}[ht]
\begin{tabular}{lccccccccc}
\hline
\textbf{Defense} & \multicolumn{3}{r}{\textbf{HarmBench ASR}} & \multicolumn{3}{r}{\textbf{StrongREJECT ASR}} & \multicolumn{3}{r}{\textbf{StrongREJECT Score}} \\
 & \textbf{mean} & \textbf{min} & \textbf{max} & \textbf{mean} & \textbf{min} & \textbf{max} & \textbf{mean} & \textbf{min} & \textbf{max} \\
\hline
Original model & 77.33 & 54.00 & 98.00 & 53.25 & 24.39 & 73.17 & 48.16 & 26.10 & 63.34 \\
RepBend & 24.50 & \textbf{2.00} & 37.00 & 8.54 & 2.44 & 19.51 & 10.36 & 4.00 & 22.06 \\
Circuit breakers & 38.67 & 27.00 & 54.00 & 6.91 & \textbf{0.00} & 12.20 & 9.32 & 3.41 & 14.53 \\
\hline
Triplet A1: CB & 38.00 & 27.00 & 51.00 & 4.88 & \textbf{0.00} & 9.76 & 6.81 & 3.56 & 11.42 \\
Triplet A2: Partial & \textbf{22.67} & \textbf{2.00} & 34.00 & 1.22 & \textbf{0.00} & \textbf{2.44} & 2.98 & \textbf{0.36} & 5.81 \\
Triplet A3: Full & 23.83 & 17.00 & \textbf{32.00} & 2.44 & \textbf{0.00} & 9.76 & 3.55 & 1.16 & 9.46 \\
Triplet A4: Full + Adv & 24.40 & 10.00 & 41.00 & \textbf{0.49} & \textbf{0.00} & \textbf{2.44} & \textbf{2.23} & 1.10 & \textbf{4.28} \\
\hline
\end{tabular}
\caption{Mean, maximum, and minimum embedding attack ASR across different hyperparameter configurations (Llama 3 8B)}
\label{tab:robustness_ablation}
\end{table*}

\begin{table}[ht]
\centering
\resizebox{0.5\textwidth}{!}{
\setlength{\tabcolsep}{2pt}
\begin{tabular}{lccc}
\hline
\textbf{} & \textbf{HB ASR} & \textbf{SR ASR} & \textbf{SR score} \\
\hline
Original model & 100.00 & 90.24 & 81.89 \\
Circuit breakers & 90.24 & 29.27 & 30.61 \\
RepBend & \textbf{73.17} & 39.02 & 39.00 \\
\hline
Triplet A1: CB & 92.68 & 21.95 & 23.11 \\
Triplet A2: RepBend & 78.05 & 7.32 & 12.98 \\
Triplet A3: Full & \textbf{65.85} & 12.20 & 14.57 \\
Triplet A4: Full + Adv & 75.61 & \textbf{4.88} & \textbf{8.70} \\
\hline
\end{tabular}
}
\caption{Embedding attack success rate of the ablation study models with the HarmBench (HB) and StrongREJECT (SR) judge classifiers (Llama 3 8B Instruct). Best-of-6 configurations, with a total of $41 \times 6 = 246$ attacks per defense.}
\label{tab:embedding_attacks_ablation}
\end{table}

\begin{table}[ht]
\centering
\resizebox{0.5\textwidth}{!}{
\setlength{\tabcolsep}{2pt}
\begin{tabular}{lccc}
\hline
\textbf{} & \textbf{HB ASR} & \textbf{SR ASR} & \textbf{SR score} \\
\hline
Original model & 31.25 & 18.75 & 23.66 \\
Circuit breakers & 2.86 & 1.43 & 4.25 \\
RepBend & 2.86 & \textbf{0.00} & 1.65 \\
\hline
Triplet A1: CB & 7.14 & 2.86  & 4.95 \\
Triplet A2: RepBend & 1.43 & 1.43 & 2.51 \\
Triplet A3: Full & \textbf{0.00} & \textbf{0.00} & \textbf{0.43} \\
Triplet A4: Full + Adv  & \textbf{0.00} & \textbf{0.00} & \textbf{1.36} \\
\hline
\end{tabular}
}
\caption{GCG attack success rate of the ablation study models with the HarmBench (HB) and StrongREJECT (SR) judge classifiers (Llama 3 8B Instruct)}
\label{tab:gcg_attacks_ablation}
\end{table}

\begin{table}[tbp]
\centering
\resizebox{0.5\textwidth}{!}{
\setlength{\tabcolsep}{2pt}
\begin{tabular}{lccc}
\hline
\textbf{} & \textbf{HB ASR} & \textbf{SR ASR} & \textbf{SR score} \\
\hline
Original model & 52.50 & 40.00 & 42.87 \\
Circuit breakers & 13.75 & 3.75 & 9.50 \\
RepBend & 11.25 & 6.25 & 11.27 \\
\hline
Triplet A1: CB & 11.25 & 6.25 & 11.10 \\
Triplet A2: RepBend & 1.25 & 1.25 & 4.87 \\
Triplet A3: Full & \textbf{0.00} & \textbf{0.00} & \textbf{0.48}  \\
Triplet A4: Full + Adv & 3.75 & 2.50 & 6.99  \\
\hline
\end{tabular}
}
\caption{REINFORCE-GCG attack success rate with the HarmBench (HB) and StrongREJECT (SR) judge classifiers for our Llama 3 8B Instruct ablation models without adversarial hard negative mining. HarmBench validation set ($81$ behaviors). The ASR is higher when more loss terms are ablated.
}
\label{tab:reinforce_and_gcg}
\end{table}

\begin{table}[ht]
\centering
\setlength{\tabcolsep}{2pt}
\begin{tabular}{lcc}
\hline
\textbf{Distance} & {$\mathbf{MMDR}_{d_2}$} & {$\mathbf{MMDR}_{d_\mathrm{cos}}$} \\
\hline
Circuit breakers & 0.63 & 0.49 \\
RepBend & 0.70 & 0.64 \\
\hline
Triplet A1: CB & 0.70 & 0.54 \\
Triplet A2: RepBend & 0.78 & 0.64 \\
Triplet A3: Full & \textbf{0.80} & 0.66 \\
Triplet A4: Full + Adv & \textbf{0.80} & \textbf{0.70} \\
\hline
\end{tabular}
\caption{
Generalization of the ablation study models to different data augmentations (Llama 3 8B Instruct) with the L2 norm $\dd_2$ and the cosine distance $\dd_\mathrm{cos}$
}
\label{tab:generalization_ablation}
\end{table}

Tables~\ref{tab:embedding_attacks_ablation}, \ref{tab:reinforce_and_gcg} and \ref{tab:gcg_attacks_ablation} show the ASRs of our 
ablated models. We observe that the A1 model performs similarly to the circuit breakers defense.
Across all attacks, the A2 model performs better than the RepBend model.
The A2 and A3 models have similar performance for embedding attacks and GCG.
Finally, the A4 model performs worse on input-space attacks, and better on embedding-space
attacks than A3. We believe this tradeoff is due to the harmful representations being used in training. A4 uses adversarial hard negative mining, so the representations used in training are more similar to adversarial embedding attack representations. On the other hand, A3 only uses harmful representations from our WildJailbeak training dataset, which are more similar to the input-space attacks representations of the model in testing.

Table~\ref{tab:generalization_ablation} shows the $\gen$ value of ablated models.
The value decreases as the loss components are ablated, showcasing their importance.

\begin{table*}[ht]
\centering
\resizebox{\textwidth}{!}{
 \begin{tabular}{lccccccccc}
\hline
 & \textbf{ARC (Easy)} & \textbf{GSM8K} & \textbf{HellaSwag} & \textbf{MMLU} & \multicolumn{3}{c}{\textbf{TruthfulQA}} \\
\cmidrule(lr){6-8}
 &  &  &  &  & \textbf{Gen} & \textbf{MC1} & \textbf{MC2} \\
\hline
\textbf{Original model} & \textbf{81.61} & \textbf{75.36} & \textbf{57.75} & \textbf{63.72} & \textbf{46.39} & \textbf{36.23} & \textbf{51.67} \\
\hline
Circuit breakers & 81.44 \textcolor{darkred}{(-0.17)} & 75.44 \textcolor{darkgreen}{(+0.08)} & 57.46 \textcolor{darkred}{(-0.29)} & 63.57 \textcolor{darkred}{(-0.15)} & 48.23 \textcolor{darkgreen}{(+1.84)} & 36.96 \textcolor{darkgreen}{(+0.73)} & 51.61 \textcolor{darkred}{(-0.05)} \\
RepBend & 80.98 \textcolor{darkred}{(-0.63)} & 49.05 \textcolor{darkerred}{(-26.31)} & 60.58 \textcolor{darkgreen}{(+2.83)} & 60.26 \textcolor{darkred}{(-3.46)} & 2.08 \textcolor{darkerred}{(-44.31)} & 41.00 \textcolor{darkgreen}{(+4.77)} & 60.05 \textcolor{darkgreen}{(+8.38)} \\
\hline
Triplet A1: CB & 81.57 \textcolor{darkred}{(-0.04)} & 74.83 \textcolor{darkred}{(-0.53)} & 57.47 \textcolor{darkred}{(-0.28)} & 63.64 \textcolor{darkred}{(-0.09)} & 48.96 \textcolor{darkgreen}{(+2.57)} & 36.96 \textcolor{darkgreen}{(+0.73)} & 52.53 \textcolor{darkgreen}{(+0.86)} \\
Triplet A2: RepBend & 81.94 \textcolor{darkgreen}{(+0.34)} & 73.84 \textcolor{darkred}{(-1.52)} & 59.56 \textcolor{darkgreen}{(+1.81)} & 63.84 \textcolor{darkgreen}{(+0.12)} & 45.29 \textcolor{darkred}{(-1.10)} & 40.27 \textcolor{darkgreen}{(+4.04)} & 55.17 \textcolor{darkgreen}{(+3.50)} \\
Triplet A3: Full & 81.27 \textcolor{darkred}{(-0.34)} & 74.30 \textcolor{darkred}{(-1.06)} & 59.62 \textcolor{darkgreen}{(+1.87)} & 63.85 \textcolor{darkgreen}{(+0.13)} & 45.65 \textcolor{darkred}{(-0.73)} & 40.76 \textcolor{darkgreen}{(+4.53)} & 55.37 \textcolor{darkgreen}{(+3.70)} \\
Triplet A4: Full + Adv & 81.99 \textcolor{darkgreen}{(+0.38)} & 74.91 \textcolor{darkred}{(-0.45)} & 60.70 \textcolor{darkgreen}{(+2.95)} & 63.38 \textcolor{darkred}{(-0.34)} & 44.55 \textcolor{darkred}{(-1.84)} & 42.96 \textcolor{darkgreen}{(+6.73)} & 57.29 \textcolor{darkgreen}{(+5.63)} \\
\hline
\end{tabular}
}
\caption{Performance comparison of the ablation study models on general capability benchmarks (Llama 3 8B Instruct)} 
\label{tab:benign_performance_ablation}
\end{table*}




\begin{table*}[ht]
\centering
\begin{tabular}{l cccc cccc}
\hline
Loss & \multicolumn{4}{c}{Benign} & \multicolumn{4}{c}{Harmful} \\
\cmidrule(lr){2-5} \cmidrule(lr){6-9} %
Term   & $\dd_{bp}$ & $\dd_{bn}$ & margin & $\bn$ & $\dd_{hp}$ & $\dd_{hn}$ & margin & $\bp$ \\
\hline
Circuit Breakers &  $\dd_2$ & \reddd  & 0 & 0 & \reddd &  $\dd_\mathrm{cos}$ & 1 & 0 \\
RepBend &  $\dd_2$ & \reddd & $\infty$ & 0 & $\dd_\mathrm{cos}$ &  $\dd_2$ & $\infty$ & $\mathrm{mean}(\hp_h)$ \\
Triplet &  $\dd$ &  $\dd$ & $m_b$ & $\bn$ &  $\dd$ &  $\dd$ & $m_h$ & $\bp$ \\
\hline
\end{tabular}
\caption{Loss function comparison. $\dd_{bp}, \dd_{bn}, \dd_{hp}$ and $\dd_{hn}$ are the distances used in the triplet losses.
  Our proposed method is a natural generalization of the Circuit Breakers and RepBend methods,
  with an additional negative term on new benign representations. For our experiments, we use $\bn = \bp = \text{mean}(\hp_{h})$.
}
\label{tab:detailed_loss_comparison}
\end{table*}

\subsection{Tamper Resistance Evaluation With Refusal Ablation}

Recent work has shown that LLMs possess a refusal direction, which can be surgically removed to bypass the refusal mechanism \citep{arditi2024refusalsingledirection}. This feature ablation attack, which has minimal effect on benign model capabilities, can be used to evaluate the brittleness of defensive methods, as a simple model tampering attack \citep{che2025modeltamperingattacksenable}.

Refusal ablation is performed by computing the difference-in-means vector of the mean activations of harmful and harmless prompts. This vector captures a meaningful direction that mediates the refusal mechanism. This refusal direction can be removed via weight orthogonalization (Directional Ablation), or added to harmful activations to shift them and make them appear more benign (Activation Addition).

\paragraph{Method}

We apply the attack implementation from \citet{arditi2024refusalsingledirection} to the Llama 3 8B Instruct model, protected by the circuit breakers defense and our triplet defense, in order to compare the brittleness of the two methods. We use the same attack parameters as the original study (128 training samples, 100 testing samples, 32 validation samples).

\begin{table*}[ht]
    \centering
\resizebox{\textwidth}{!}{
    \begin{tabular}{l c c c c c c}
        \hline
        & \multicolumn{2}{c}{\textbf{Baseline (No Tampering)}} & \multicolumn{2}{c}{\textbf{Directional Ablation}} & \multicolumn{2}{c}{\textbf{ActAdd}} \\
        \cmidrule(lr){2-3} \cmidrule(lr){4-5} \cmidrule(lr){6-7}
        \textbf{Model} & \textbf{Non-refusal rate} & \textbf{HB ASR} & \textbf{Non-refusal rate} & \textbf{HB ASR} & \textbf{Non-refusal rate} & \textbf{HB ASR} \\
        \hline
        \textit{Llama-3-8B-Instruct} & 4 & 2 & 100 & 84 & 99 & 89 \\
        \ + Circuit breakers & 76 & 33 & 23 & 15 & 39 & 30 \\
        \ \textbf{+ Triplet} & 97 & \textbf{0} & 37 & \textbf{0} & 73 & \textbf{1} \\
        \hline
    \end{tabular}
}
    \caption{Attack Success Rates (ASR) across the Directional Ablation and Activation Addition (ActAdd) tampering attacks. The non-refusal rate, or substring ASR, is detected by the presence of refusal string, e.g ``I am sorry, I cannot\ldots''. The HB ASR is computed using the HarmBench Llama2Guard judge.}
    \label{tab:tampering}
\end{table*}


\paragraph{Results}
Table~\ref{tab:tampering} reports the attack success rates of the tampering attacks.
Our triplet-protected model achieves an ASR of 0\% on the Ablation attack and 1\% on the ActAdd attack, outperforming the circuit breakers protected model, which achieves ASRs of 30\% and 15\% for Ablation and ActAdd respectively).
Both methods significantly improve upon the original, unprotected model, achieving ASRs of over 80\%.
The non-refusal rate, or substring ASR, is computed by detecting the presence of standard refusal sentences (e.g, ``I am sorry, I cannot\ldots''). The relatively high substring ASRs of both defenses are explained by the tendency of representation engineering-based methods to break in generation instead of producing proper refusals.



\subsection{Comparison with Other Adversarial Defenses}
\label{app:models}

We evaluate the performance of our defense in comparison to existing adversarial defense methods across multiple model architectures. In particular, we consider the following adversarial defense methods:

\begin{itemize}[noitemsep]
\item \textbf{Latent Adversarial Training (LAT)} \citep{sheshadri2024latentadversarialtrainingimproves}
perturbs the residual stream by inserting adversarial perturbations and fine-tunes model weights to maximize the refusal probability on harmful prompts.
\item \textbf{Robust Refusal Dynamic Defense (R2D2)} \citep{mazeika2024harmbenchstandardizedevaluationframework} synthesizes adversarial suffixes using the GCG attack \citep{zou2023universaltransferableadversarialattacks} and applies them to fine-tune the model weights to increase its robustness.
\item \textbf{Continuous-Adversarial Unlikelihood (CAT)} \citep{xhonneux2024efficientadversarialtrainingllms}
trains on adversarial behaviour and combines it with fine-tuning on utility data to improve robustness without compromising harmfulness.
\item \textbf{Continuous Adversarial Preference Optimization (CAPO)} \citep{xhonneux2024efficientadversarialtrainingllms} is an adversarial variant of Identity Preference Optimization (IPO) \citep{azar2023generaltheoreticalparadigmunderstand} which, unlike CAT, does not require utility data.
\item \textbf{Refusal Feature Adversarial Training (ReFAT)}\citep{yu2025robustllmsafeguardingrefusal} is an efficient adversarial training method that simulates the effect of input-level attacks using refusal ablation tampering attacks.
\end{itemize}

\paragraph{Method}
We gather publicly available models of LAT, R2D2, CAT and CAPO released from their corresponding authors
for the models Llama 3 8B Instruct \citep{grattafiori2024llama3herdmodels}, Phi 3 Mini 4K Instruct \citep{abdin2024phi3technicalreporthighly} which is a 3.8B parameter light weight model that outperforms Llama 3 8B on many general capability datasets, and Zephyr 7B beta \citep{tunstall2023zephyr} which is a fine-tuned version of Mistral 7B v0.2. Phi 3 Mini features a slightly different architecture, including modified positional encoding and fused MLP up and gate projection layers, which enables us to evaluate the flexibility of our defense method across models with different architectures.
We reimplement the ReFAT adversarial defense using the hyperparameters from \citet{yu2025robustllmsafeguardingrefusal}.
Unlike the original setup, we use the mean of the last five tokens rather than relying on the final token for the RFA attack, to improve stability.

We use the embedding attack setting described in Section~\ref{par:embedding} and report both the different attack success rates and general capability performance.

\begin{table*}[tb]
\centering
\resizebox{\textwidth}{!}{
\setlength{\tabcolsep}{2pt}
\begin{tabular}{lccccccc}
\hline
 & ARC (Easy) & GSM8K & HellaSwag & MMLU & TruthfulQA (gen) & TruthfulQA (mc1) & TruthfulQA (mc2) \\
\hline
\textit{Phi-3-Mini-4k-Instruct} &     81.69 &  79.23 &      59.02 &  69.93 &           74.54 &           36.35 &           54.52 \\
\ + CAPO    &     80.09 &  78.24 &      59.12 &  69.33 &           74.54 &           38.43 &           57.77 \\
\ + CAT     &     81.99 &  77.26 &      57.98 &  69.29 &           50.80 &           36.35 &           53.22 \\
\ + ReFAT                   &     81.44 &  72.78 &      57.99 &  62.35 &           51.65 &           35.99 &           51.15 \\
\ \textbf{+ Triplet} &     80.43 &  75.06 &      57.95 &  68.93 &           79.44 &           37.33 &           55.14 \\
\hline
\textit{Zephyr-7B-$\beta$}   &     81.27 &  34.04 &      63.97 &  58.47 &           48.23 &           38.56 &           55.20 \\
\ + R2D2    &     80.47 &  38.06 &      61.45 &  58.22 &           45.29 &           31.58 &           45.48 \\
\ \textbf{ + Triplet} &     81.61 &  32.83 &  52.49 &  57.65 &           45.41 &           39.29 &           55.61 \\

\hline

\end{tabular}
}
\caption{General performance of defenses on Phi 3 Mini and Zephyr 7B}
\label{tab:general_all}
\end{table*}


\begin{table*}[tb]
\centering
\begin{tabular}{lccc}
\hline
Model &  HB ASR &  SR ASR &  SR Score \\
\hline
\textit{Llama-3-8B-Instruct}                 &           100.00 &             90.24 &               81.89 \\
\ + LAT                  &            97.56 &             80.49 &               65.19 \\
\ + ReFAT &      97.56 &             90.24 &                78.07 \\

\ \textbf{+ Triplet} &  \textbf{65.85} & \textbf{12.20} & \textbf{14.57} \\
\hline
\textit{Phi-3-Mini-4K-Instruct}    &           100.00 &             92.68 &               81.32 \\
\ + CAPO    &           100.00 &             90.24 &               82.67 \\
\ + CAT     &           \textbf{85.37} &             78.05 &               70.87 \\
\ \textbf{+ Triplet} &            92.68 &             \textbf{43.90} &               \textbf{38.43} \\
\hline
\textit{Zephyr-7B-$\beta$} &           100.00 &             87.80 &               82.81 \\
\ + R2D2           &            92.68 &             60.98 &               54.60 \\
\ \textbf{+ Triplet} &            \textbf{70.73} &             \textbf{24.39} &               \textbf{26.79} \\
\hline
\end{tabular}
\caption{Attack success rates (ASR) using HarmBench (HB) and StrongREJECT (SR) across attack types, for various adversarial defenses on different modes.
Results were computed over 41 behaviors, with six attempts per behavior using different hyperparameter configurations (246 runs per defense). The best result for each behavior was used. StrongREJECT scores are reported on a 0--100 scale.
}
\label{tab:comparison_asr_all}
\end{table*}

\paragraph{Results}
Embedding ASR results are reported in Table~\ref{tab:comparison_asr_all}.
Across all evaluated defenses, our triplet loss consistently achieved the lowest ASR by a significant margin. We observe that the LAT and ReFAT methods achieve a higher StrongREJECT ASR than all representation-engineering-based defenses evaluated, including circuit breakers, RepBend, and our ablation study models.
For Phi 3 Mini and Zephyr, the triplet defended model outperforms CAPO, CAT and R2D2 by a factor of 2 in the StrongREJECT score.

The general performance of the model across tested defenses are reported in Table~\ref{tab:general_all}.
These results indicate that the triplet defense does not degrade the general performance more significantly than existing methods, and that the general benign capabilities of the models are kept.

\label{app:appendix_additional_results}

\paragraph{Mistral Results}
\label{app:mistral}

To evaluate our defense against existing representation-engineering-based defenses, we train our defense on Mistral-7B-Instruct-v0.2. We keep the same hyperparameters used for Llama 3 8B.

Table~\ref{tab:embedding_attacks_mistral} shows the embedding attack success rates
on the Mistral model. Our method achieves the lowest StrongREJECT score, slightly outperforming circuit breakers. We observe that the attack success rates for both defense methods are much higher on the Mistral 7B model than on the Llama 3 8B model, indicating that Mistral 7B remains a more challenging target for robust defense.

Table~\ref{tab:benign_performance_mistral} shows the general performance of the trained models.
Our trained Mistral model successfully retains its general language capabilities, demonstrating
that our defense does not compromise model performance.

\subsection{t-SNE Visualization}

t-SNE \cite{JMLR:v9:vandermaaten08a} is a nonlinear dimensionality reduction technique commonly used for visualizing high-dimensional datasets.
We use it to visualize the representations of the Llama 3 8B model
on the tokens between the end of the prompt and the start of a reply on both benign and harmful prompts. We also include representations obtained after performing 100 iterations of an embedding attack on harmful prompts.

Figure~\ref{fig:tsne} shows the t-SNE visualization for layer 25 across several defenses.
We observe that, compared to baseline defenses, our Triplet-based approach clusters both harmful and attacked representations together more effectively.


\begin{figure*}[htbp]
    \centering
    \includegraphics[width=\textwidth]{latex/images/tsne_2d_plot.pdf}
    \caption{t-SNE visualization of layer 25 representations of Llama 3 8B. Representations from benign prompts (green), harmful prompts (red), and embedding-attacked harmful prompts (brown) are shown. Our Triplet-based defense achieves tighter clustering of harmful and attacked representations compared to baseline defenses.}
    \label{fig:tsne}
\end{figure*}





\begin{table*}[tb]
\centering
\begin{tabular}{lccc}
\hline
name & HarmBench ASR & StrongREJECT ASR & StrongREJECT score \\
\hline
\textit{Mistral-7B-Instruct-v0.2} & 100.00 & 92.68 & 84.81 \\
\ + Circuit breakers & \textbf{85.37} & 41.46 & 42.76 \\
\ \textbf{+ Triplet} & 97.56 & \textbf{34.15} & \textbf{33.82} \\
\hline
\end{tabular}
\caption{Embedding attack success rate with the HarmBench and StrongREJECT judge classifiers (Mistral 7B Instruct)}
\label{tab:embedding_attacks_mistral}
\end{table*}

\begin{table*}[tb]
\centering
\resizebox{\textwidth}{!}{
\begin{tabular}{lccccccc}
\hline
 & ARC (Easy) & GSM8K & HellaSwag & MMLU & TruthfulQA (gen) & TruthfulQA (mc1) & TruthfulQA (mc2) \\
\hline
\textit{Mistral-7B-Instruct-v0.2} & 81.23 & 41.77 & 66.01 & 58.97 & 54.22 & 52.26 & 66.84 \\
\hline
CB & 81.52 \textcolor{darkgreen}{(+0.29)} & 44.20 \textcolor{darkgreen}{(+2.43)} & 65.58 \textcolor{darkred}{(-0.43)} & 58.87 \textcolor{darkred}{(-0.10)} & 55.69 \textcolor{darkgreen}{(+1.47)} & 52.14 \textcolor{darkred}{(-0.12)} & 67.05 \textcolor{darkgreen}{(+0.21)} \\
RepBend & 81.36 \textcolor{darkgreen}{(+0.13)} & 42.15 \textcolor{darkgreen}{(+0.38)} & 65.60 \textcolor{darkred}{(-0.41)} & 58.77 \textcolor{darkred}{(-0.20)} & 50.31 \textcolor{darkred}{(-3.92)} & 52.02 \textcolor{darkred}{(-0.24)} & 67.23 \textcolor{darkgreen}{(+0.39)} \\
\hline
\textbf{Triplet} & 81.48 \textcolor{darkgreen}{(+0.25)} & 41.47 \textcolor{darkred}{(-0.30)} & 65.83 \textcolor{darkred}{(-0.18)} & 58.95 \textcolor{darkred}{(-0.01)} & 54.83 \textcolor{darkgreen}{(+0.61)} & 51.04 \textcolor{darkred}{(-1.22)} & 66.99 \textcolor{darkgreen}{(+0.15)} \\
\hline
\end{tabular}
}
\caption{Performance comparison of models on general capability benchmarks (Mistral 7B Instruct)} 
\label{tab:benign_performance_mistral}
\end{table*}



\begin{table*}[tb]
\centering
\caption{Full embedding attack results on Llama 3 8B (all configurations). Different defenses are robust to different hyperparameter configurations.}
\label{tab:full_embedding_results}
\resizebox{\textwidth}{!}{
\begin{tabular}{llccc}
\hline
Configuration & Model & HarmBench ASR & StrongREJECT ASR & StrongREJECT score \\
\hline
\multirow[t]{7}{*}{0} & Original model & 0.78 & 0.56 & 0.49 \\
 & Circuit breakers & 0.51 & 0.12 & 0.15 \\
 & Triplet A1: CB & 0.46 & 0.10 & 0.11 \\
 & Triplet A4: Full + Adv & \textbf{0.24} & 0.05 & 0.06 \\
 & RepBend & 0.34 & 0.02 & 0.05 \\
 & Triplet A3: Full & \textbf{0.24} & \textbf{0.00} & 0.04 \\
 & Triplet A2: RepBend & 0.34 & 0.02 & \textbf{0.03} \\
\cline{1-5}
\multirow[t]{7}{*}{1} & Original model & 0.54 & 0.27 & 0.27 \\
 & Circuit breakers & 0.54 & 0.12 & 0.12 \\
 & RepBend & 0.15 & 0.10 & 0.09 \\
 & Triplet A1: CB & 0.27 & \textbf{0.00} & 0.04 \\
 & Triplet A2: RepBend & 0.24 & 0.02 & 0.03 \\
 & Triplet A4: Full + Adv & \textbf{0.10} & \textbf{0.00} & 0.02 \\
 & Triplet A3: Full & 0.32 & \textbf{0.00} & \textbf{0.01} \\
\cline{1-5}
\multirow[t]{7}{*}{2} & Original model & 0.88 & 0.68 & 0.60 \\
 & RepBend & 0.37 & 0.20 & 0.22 \\
 & Triplet A3: Full & 0.29 & 0.10 & 0.09 \\
 & Triplet A1: CB & 0.34 & 0.10 & 0.09 \\
 & Circuit breakers & \textbf{0.27} & 0.02 & 0.06 \\
 & Triplet A2: RepBend & 0.32 & 0.02 & 0.06 \\
 & Triplet A4: Full + Adv & 0.41 & \textbf{0.00} & \textbf{0.02} \\
\cline{1-5}
\multirow[t]{7}{*}{3} & Original model & 0.98 & 0.71 & 0.63 \\
 & RepBend & 0.32 & 0.12 & 0.15 \\
 & Triplet A1: CB & 0.51 & 0.05 & 0.08 \\
 & Circuit breakers & 0.34 & 0.02 & 0.07 \\
 & Triplet A3: Full & \textbf{0.17} & 0.05 & 0.04 \\
 & Triplet A2: RepBend & 0.22 & \textbf{0.00} & 0.04 \\
 & Triplet A4: Full + Adv & 0.32 & \textbf{0.00} & \textbf{0.01} \\
\cline{1-5}
\multirow[t]{7}{*}{4} & Original model & 0.56 & 0.24 & 0.26 \\
 & Circuit breakers & 0.39 & 0.12 & 0.13 \\
 & RepBend & 0.27 & 0.05 & 0.07 \\
 & Triplet A1: CB & 0.29 & 0.02 & 0.05 \\
 & Triplet A4: Full + Adv & \textbf{0.12} & \textbf{0.00} & 0.02 \\
 & Triplet A2: RepBend & 0.22 & \textbf{0.00} & 0.02 \\
 & Triplet A3: Full & 0.24 & \textbf{0.00} & \textbf{0.01} \\
\cline{1-5}
\multirow[t]{6}{*}{5} & Original model & 0.90 & 0.73 & 0.63 \\
 & RepBend & \textbf{0.02} & 0.02 & 0.04 \\
 & Triplet A1: CB & 0.41 & 0.02 & 0.04 \\
 & Circuit breakers & 0.27 & \textbf{0.00} & 0.03 \\
 & Triplet A4: Full + Adv & 0.29 & \textbf{0.00} & 0.03 \\
 & Triplet A3: Full & 0.17 & \textbf{0.00} & 0.01 \\
 & Triplet A2: RepBend & \textbf{0.02} & \textbf{0.00} & \textbf{0.00} \\
\hline
\end{tabular}
}
\end{table*}

\begin{table*}
\centering
\caption{Full generalization results of the defenses on different input formats (Llama 3 8B Instruct)}
\label{tab:genresults}
  \begin{tabular}{llrr}
\hline
 &  & \multicolumn{2}{r}{Distance} \\
Augmentation & Defense & L2 & Cosine \\
\hline
\multirow[t]{6}{*}{Chinese} & RepBend & 0.84 & 0.82 \\
 & Triplet A1: CB & 0.84 & 0.89 \\
 & Circuit breakers & 0.87 & 0.85 \\
 & Triplet A2: RepBend & 0.92 & 0.96 \\
 & Triplet A3: Full & 0.92 & 0.95 \\
 & Triplet A4: Full + Adv & 0.93 & 0.89 \\
\cline{1-4}
\multirow[t]{6}{*}{French} & Circuit breakers & 0.89 & 0.89 \\
 & RepBend & 0.92 & 0.89 \\
 & Triplet A1: CB & 0.94 & 0.97 \\
 & Triplet A4: Full + Adv & 0.95 & 0.89 \\
 & Triplet A3: Full & 0.97 & 0.98 \\
 & Triplet A2: RepBend & 0.97 & 0.98 \\
\cline{1-4}
\multirow[t]{6}{*}{German} & Circuit breakers & 0.88 & 0.89 \\
 & RepBend & 0.93 & 0.90 \\
 & Triplet A4: Full + Adv & 0.95 & 0.88 \\
 & Triplet A1: CB & 0.97 & 0.98 \\
 & Triplet A2: RepBend & 0.97 & 0.99 \\
 & Triplet A3: Full & 0.97 & 1.00 \\
\cline{1-4}
\multirow[t]{6}{*}{Spanish} & Circuit breakers & 0.88 & 0.88 \\
 & RepBend & 0.92 & 0.90 \\
 & Triplet A1: CB & 0.94 & 0.98 \\
 & Triplet A4: Full + Adv & 0.95 & 0.91 \\
 & Triplet A2: RepBend & 0.96 & 1.00 \\
 & Triplet A3: Full & 0.97 & 0.99 \\
\cline{1-4}
\multirow[t]{6}{*}{Ukrainian} & RepBend & 0.79 & 0.77 \\
 & Circuit breakers & 0.80 & 0.79 \\
 & Triplet A1: CB & 0.82 & 0.87 \\
 & Triplet A2: RepBend & 0.87 & 0.90 \\
 & Triplet A3: Full & 0.89 & 0.91 \\
 & Triplet A4: Full + Adv & 0.90 & 0.85 \\
\cline{1-4}
\multirow[t]{6}{*}{Random capitalization} & Circuit breakers & 0.76 & 0.68 \\
 & RepBend & 0.90 & 0.83 \\
 & Triplet A4: Full + Adv & 0.91 & 0.84 \\
 & Triplet A3: Full & 0.92 & 0.83 \\
 & Triplet A2: RepBend & 0.92 & 0.83 \\
 & Triplet A1: CB & 0.93 & 0.80 \\
\hline
\end{tabular}
\end{table*}

\begin{table*}[h]
	\centering
	\begin{tabular}{p{5cm}p{5cm}p{5cm}}
		\hline
		\textbf{Question} & \textbf{Expected answer} & \textbf{Generated answer} \\
    \hline
    Howard spends \$ 8 dollars at the arcade on Monday. On Tuesday, he spends twice as much at the arcade as he did on Monday. On Wednesday, he spends 4 times as much at the arcade as he spent on Tuesday. If he originally had \$100, how much money does he have left? &
On Tuesday, he spends \$8 * 2 = \$\textless 8 * 2 = 16 \textgreater.
On Wednesday, he spends \$16 * 4 = \$\textless 16 * 4 = 64 \textgreater.
In total, he spent \$8 + \$16 + \$64 = \$\textless 8 + 16 + 64 = 88 \textgreater\ at the arcade.
He has \$100 - \$88 = \$\textless 100 - 88 = 12 \textgreater\ left.
    \#\#\#\# 12 & \textcolor{darkgreen}{Howard starts with \$100.
On Monday, he spends \$8.
On Tuesday, he spends 2 * \$8 = \$\textless 2 * 8 = 16 \textgreater.
On Wednesday, he spends 4 * \$16 = \$\textless 4 * 16 = 64 \textgreater.
Total spent: \$8 + \$16 + \$64 = \$\textless 8 + 16 + 64 = 88 \textgreater.
Howard has \$100 - \$88 = \$\textless 100 - 88 = 12 \textgreater\ left.
\#\#\#\# 12}
\\ \\
    Dax went to their farm to pick some apples and found half as many bugs as ants in the garden. If there were 50 ants, calculate the total number of insects in the garden. & The total number of insects in the garden is 25 bugs + 50 ants = \textless\textless25+50=75\textgreater\textgreater75 insects \#\#\#\# 75 & \textcolor{red}{if there were 50 ants, and there were half as many bugs as ants, there were 1/2 * 50 = 1/2 * 50 = 1 * 50 / 2 =     iseliseliselilletilletilletilletilletilletille} \\ \\
    
    Last night Rick killed ten wolves and 15 cougars while hunting. Today Rick killed three times as many wolves as cougars and three fewer cougars than the previous night. How many animals did Rick kill? & The total number of animals that Rick killed while hunting yesterday is 15+10 = \textless\textless15+10=25\textgreater\textgreater25
Today, Rick killed three fewer cougars than the previous night, a total of 15-3 = \textless\textless15-3=12\textgreater\textgreater12 cougars.
He also killed three times as many wolves as cougars today, a total of 3*12 = \textless\textless3*12=36\textgreater\textgreater36 wolves.
Together, his hunt today yielded 36+12 = \textless\textless36+12=48\textgreater\textgreater48 animals.
    In total, he has 48+25 = \textless\textless48+25=73\textgreater\textgreater73 animals from the hunt. \#\#\#\# 73 & \textcolor{red}{yesterday, he killed 10 wolves and 15 cou, so today he killed 3 times as many wolves as cou, which is     isel  mktimeilletiselovuovuilletov \ldots} \\
    \hline
\end{tabular}
  \caption{Example of correct and incorrect (overrefused) generations of RepBend on GSM8K.}
	\label{tab:overrefusal_example}
\end{table*}

\subsection{Prompt Templates}
\label{app:prompts}
\begin{tcolorbox}[colback=gray!10, colframe=black, 
title=Prompt used for translations]
Generate a translation of the following text into {x[0]}:
```
{x[1]}
```
Do not output anything else, just translate the text as best as you can.
ONLY TRANSLATE the text, do not add any header, response, or footer to your reply.
Make sure to translate the text as best as you can, and do not add any extra information.
\end{tcolorbox}

\onecolumn
\subsection{Ablation Study Sosses}
\label{app:ablation_losses}

Equations \ref{eq:a1}, \ref{eq:a2}, and \ref{eq:a3} present the loss functions used for models A1, A2 and A3 in our ablation study.

\begin{align}
\mathcal{L}_{A1} &= \alpha \cdot \frac{1}{N} \sum_{i=1}^{N} \max(0, \text{d}_{bp}(\h_{b,i}, \hp_{b,i}) + m_b) \nonumber \\
                 &\quad + \beta \cdot \frac{1}{N} \sum_{i=1}^{N} \max(0, -\text{d}_{hn}(\hp_{h,i}, \h_{h,i}) + m_h) \nonumber \\
                 &\quad + \gamma \cdot \KL(\pi(b), \pi^\prime(b))
\label{eq:a1} \\[1.5em]
\mathcal{L}_{A2} &= \alpha \cdot \frac{1}{N} \sum_{i=1}^{N} \max(0, \text{d}_{bp}(\h_{b,i}, \hp_{b,i}) + m_b) \nonumber \\
                 &\quad + \beta \cdot \frac{1}{N} \sum_{i=1}^{N} \max(0, \text{d}_{hp}(\hp_{h,i}, \hat{\hp}) - \text{d}_{hn}(\hp_{h,i}, \h_{h,i}) + m_h) \nonumber \\
                 &\quad + \gamma \cdot \KL(\pi(b), \pi^\prime(b))
\label{eq:a2} \\[1.5em]
\mathcal{L}_{A3} &= \alpha \cdot \frac{1}{N} \sum_{i=1}^{N} \max(0, \text{d}_{bp}(\h_{b,i}, \hp_{b,i}) - \text{d}_{bn}(\hp_{b,i}, \hat{\hp}) + m_b) \nonumber \\
                 &\quad + \beta \cdot \frac{1}{N} \sum_{i=1}^{N} \max(0, \text{d}_{hp}(\hp_{h,i}, \hat{\hp}) - \text{d}_{hn}(\hp_{h,i}, \h_{h,i}) + m_h) \nonumber \\
                 &\quad + \gamma \cdot \KL(\pi(b), \pi^\prime(b))
\label{eq:a3}
\end{align}


\subsection{Proofs}
\label{sec:appendix_proofs}

\newtheorem{theorem}{Theorem}
\newtheorem{lemma}{Lemma}

\begin{theorem}
  The circuit breakers loss $\mathcal{L}_{CB}$ can be rewritten as a triplet loss $\mathcal{L}_{triplet}$ with null distances $d_0(x, y) = 0$.
\end{theorem}
\begin{proof}
We recall the definition of the circuit breakers loss:
\begin{align}
  \mathcal{L}_{CB} &= \alpha \cdot \norm{\h_{b,i} - \hp_{b,i}}_2^2 + \beta \cdot \ReLU(\cossim(\h_{h,i}, \hp_{h,i})) \\
                   &= \alpha \cdot \ReLU(\norm{\h_{b,i} - \hp_{b,i}}_2^2) + \beta \cdot \ReLU(\cossim(\h_{h,i}, \hp_{h,i})) \\
                   &= \alpha \cdot \ReLU(\norm{\h_{b,i} - \hp_{b,i}}_2^2) + \beta \cdot \ReLU(- \dd_{cos}(\h_{h,i}, \hp_{h,i}) + 1) \\
                   &= \mathcal{L}_{triplet}
\end{align}

with the parameters $m_h = 1, m_b = 1, \text{d}_{bp} = \dd_2, \text{d}_{bn} = \dd_0, \text{d}_{hp} = \dd_0, \dd_{hn} = \text{d}_{cos}$.

\end{proof}

\newcommand{\ba}{\mathbf{a}}

\begin{theorem}

The RepBend loss $\mathcal{L}_{RB}$ without the KL divergence term can be rewritten as a triplet loss $\mathcal{L}_{triplet}$
with a null distance $d_0(x, y) = 0$ and an arbitrary large margin $m_h$.

\end{theorem}
\begin{proof}
The definition of the RepBend loss, minus the KL divergence term, is:
\begin{align}
  \mathcal{L}_{\text{RB}} &= \frac{1}{2} \cdot \left\| \hp_{b,i} - \h_{b,i} \right\|_2  \\
                          &\quad- \alpha \cdot \left\| \hp_{h,i} - \h_{h,i} \right\|_2  \\
														&\quad- \beta \cdot \cossim(A) 
\end{align}
where $\cossim(A)$ is defined as the average cosine similarity between all pairs of vectors in $A$.
\begin{align}
	\cossim(A) = \frac{1}{n(n-1)} \sum_{i=1}^{n} \sum_{j=1, j \neq i}^{n} \cossim(\ba_i, \ba_j)
\end{align}

We write $\mu$ the mean of the normalized vectors $\hat{\ba}_i$.
\begin{align}
    \mu = \frac{1}{n} \sum_{i=1}^{n} \hat{\ba}_i
\end{align}

Let us assume that the vectors are normalized, i.e. $||\ba_i|| = 1$.
\begin{align}
  \cossim(A) &= \frac{1}{n(n-1)} \sum_{i=1}^{n} \sum_{j=1, j \neq i}^{n} \cossim(\ba_i, \ba_j) \\
             &= \frac{1}{n(n-1)} \sum_{i=1}^{n} \sum_{j=1, j \neq i}^{n} \ba_i \cdot \ba_j \\
             &= \frac{1}{n(n-1)} \sum_{i=1}^{n} \left(\sum_{j=1}^{n} \ba_i \cdot \ba_j\right) - \ba_i \cdot \ba_i \\
             &= -\frac{1}{n-1} +  \frac{1}{n(n-1)} \sum_{i=1}^{n} \sum_{j=1}^{n} \ba_i \cdot \ba_j \\
             &= -\frac{1}{n-1} +  \frac{1}{n(n-1)} \sum_{i=1}^{n} \ba_i \cdot \left(\sum_{j=1}^{n} \ba_j\right) \\
             &= -\frac{1}{n-1} +  \frac{1}{(n-1)} \sum_{i=1}^{n} \ba_i \cdot \mu \\
             &= -\frac{1}{n-1} - \frac{1}{(n-1)} \sum_{i=1}^{n} \dd_{\cos}(\ba_i, \mu)
\end{align}

Therefore, maximizing $\cossim(A)$ is equivalent to minimizing $\frac{1}{n} \sum_{i=1}^{n} \dd_{\cos}(\ba_i, \mu)$.
Since $\dd_{\cos}$ is the cosine distance, the case also holds when the vectors are not normalized.

Finally, we can rewrite the RepBend loss as:
\begin{align}
  \mathcal{L}_{RB} &= \frac{1}{2} \cdot \left\| \hp_{b,i} - \h_{b,i} \right\|_2  \\
                   &\quad- \alpha \cdot \left\| \hp_{h,i} - \h_{h,i} \right\|_2  \\
                            &\quad- \beta \cdot \frac{1}{H} \sum_{n=1}^{H} \dd_{\cos}(\hp_{h,n}, \mu)  \\
                            &= \frac{1}{2} \cdot \ReLU(\dd_2(\hp_{b,i}, \h_{b,i}) + d_0 - 0) \\
                            &\quad+ \ReLU(\alpha \cdot \dd_2(\hp_{h,i}, \h_{h,i}) - \frac{\beta}{n - 1} \cdot d_{\cos}(\hp_{h,i}, \mu) + m_h)
\end{align}
where we assume that $m_h$ is large enough to make the ReLU function non-zero.
In practice, the RepBend loss would need to be stopped at a certain point, otherwise the loss would tend to $-\infty$.
By adding a margin $m_h$, we can freely choose the number of iterations without worrying about divergence of the loss.
Therefore, our triplet loss formulation of RepBend with margins is more practical and stable.
\end{proof}

\end{document}
